\chapter{消耗战中的资源陷阱}


毋庸讳言，第五次反“围剿”开始前后，中央苏区面临着空前的困难，困难既来自外在的压迫，同时也有内生瓶颈的制约。由于国民党军的全面封锁，中央苏区活动空间不断被压缩，人力、物质、经济、政治资源都日显匮乏，中共在苏区的生存面临着越来越多的挑战，作战资源成为中共在即将到来的消耗战中的最大难题。

\section{结构性的限制——中央苏区的人力与物力}

\subsection{人力资源}


由于国民党军在第五次“围剿”中采取持久作战方针，而中共面对国民党军的堡垒策略和压迫式战法，选择了留在苏区的应战之路，这意味着一场大规模的持久消耗战不可避免。对于一场双方在战略上已没有多少悬念的战争而言，战争资源对双方胜败有着关键性的影响。虽然，中央苏区经过数年的发展，已拥有一定规模，但当国民党开始从战争资源上寻求突破口时，中共占据的一隅之地和有限资源仍然无法真正与南京政府抗衡。中共自己在当时就承认：“照物质上的力量比较起来，白军真可以两三个月‘荡平赤匪’。”\footnote{狄康：《庐山会议的大阴谋》，上海中央局：《斗争》第49～50期，1933年8月。}而其高级干部在日后也谈道：“我们中央苏区最后不能坚持而被迫万里长征，除了路线错误之外，其中重要原因之一就是我们财政的枯竭，人力物力财力的枯竭，几乎到了山穷水尽的境地。”\footnote{舒同：《贯彻大会的精神与方针——舒主任在高干会的总结报告》，《斗争》第4期，1947年11月30日。}“根据地人民竭尽全力，也不能保证及时供应，这也是导致第五次反‘围剿’失败的原因之一。”\footnote{蔡长风：《征途漫忆》，海潮出版社，1994，第12页。}


在消耗作战中，人力资源是决定战争成败的重要一环。作为全国广大区域控制者，南京政府充足的人力资源可保证其兵力来源无忧，而中央苏区地域有限，且位于人口稀少的山地地区，数量上处于绝对劣势。


中央苏区的人口，没有准确的统计数据。1932年上半年，由于日本发动九一八事变，南京政府“剿共”军事全面收缩，中央苏区空前发展，据此时中共方面的报告，江西苏区“居民有二百四十五万以上”。\footnote{《江西苏区中共省委工作总结报告（1932年5月）》，《中央革命根据地史料选编》（上），第425页。}不过这只是一个粗略的估计。从区域角度看，当时整个赣南区域（即今赣州市辖区）人口为200多万，属于中央苏区者占面积一半左右，估计人口在百万以上。同时苏区面积还延伸到赣中、赣东一些地区，这里有数十万人。两者相加，江西苏区人口应在200万人以上。如果分县细加衡量，根据可以搜集到的材料看，江西几个主要苏区县宁都、瑞金、广昌三县苏维埃革命前人口分别为322480、310000、110000人，\footnote{《江西各县人口变动表》，汪浩：《收复匪区之土地问题》，正中书局，1935，第44页。}会昌县215406人，\footnote{南京政府内政部调查报告。}这四个县共有95万多人。但应该注意的是，苏维埃革命开始后，这里的人员变化很大，1932年5月的统计，宁都为204651人，瑞金、会昌均为24万余人，\footnote{《县区乡人口统计表》，《中央革命根据地史料选编》上册，第454页。}宁都、瑞金均大幅减少，会昌有所增加。红军长征后的统计分别为161240、200000、68000、154404人，\footnote{《江西各县人口变动表》，汪浩：《收复匪区之土地问题》，第44页。}仅相当于革命前的50\%、64.51\%、61.82\%、71.7\%。以长征后总人数除苏维埃革命前总人数，三县长征后总人数只相当于革命前的60.9\%。这些伤亡有些是在第五次反“围剿”期间产生的，第五次反“围剿”前如以80\%计，上述四县人数为76万余人。1930年代石城全县人口13.6万余人。\footnote{温昌义：《石城红军知多少》，《石城文史资料》第2辑，政协石城文史资料委员会编印，1987，第55页。}于都境内并存着两个县，胜利县人口15.33余万，于都县人口19.1余万。\footnote{刘熙朋、邹卫东：《苏区于都红色政权组织发展史略》，《于都文史资料》第2辑，政协于都文史资料委员会编印，1991，第4页。}赣县苏区人口近20万。万泰8万人。这几县相加也为76万。再加上信丰（1932年底的报告苏区人口为31300多人\footnote{《中共信丰县委两个月（十月二十日至十二月二十日）冲锋工作报告》，《江西革命历史文件汇集（1932年）》（二），第460页。}）、永丰（苏区、半苏区人口110200余人\footnote{《中共永丰县委两个月（十月二十日至十二月二十日）冲锋工作报告》，《江西革命历史文件汇集（1932年）》（二），第457页。}）、安远（1928年安远县共有人口175720人，1935年安远人口116062人\footnote{《编组保甲报告》，谢才丰：《旧时安远的户口》，政协安远文史资料委员会编印，1989，第58页。}）、宜黄、寻乌、南丰、黎川等部分苏维埃化的县份，将第五次反“围剿”开始前后江西苏区人口定位于200万人以上应属适中。


在江西苏区空前扩大同时，1932年由于十九路军入闽，红军在福建采取收缩战略，闽西苏区面积有所缩小。1933年6月，红军组成东征军入闽，改变十九路军入闽以来红军在闽西的防御和被压缩趋势，苏区面积又有扩大。闽西苏区基本区域为长汀、连城、永定、建宁、泰宁、宁化、上杭、武平等数县。迄未发现苏维埃时期详细的人口统计资料。根据1940年的统计，闽西数县人口是：长汀195000人，连城104999人，宁化121488人，武平152471人，上杭193837人，龙岩136930人，清流56143人，永定166714人，明溪34942人，建宁56856人，泰宁47737人，这些县人口总计1267097人。\footnote{《福建省各县区农业状况》上、下册，福建省农林处统计室编印，1942，本书据其所载各县区农业概况中的人口数据统计。}据1947年的统计，闽西数县人口是：长汀198200人，上杭196188人，永定171486人，龙岩141632人，武平141125人。\footnote{《福建省各县乡镇区户口统计表》，福建省档案馆藏档435011-6-3846。}两个数据比较，相同的县份人口仅有缓慢增长，1940年代初的数据虽可能比之苏维埃时期略高，但还是有相当的参考价值。上述县份除长汀、上杭等外，辖境不完全属于苏区范围，另有部分苏区地域不在上述县份。出入相抵，闽西苏区总人口当在百万左右。闽西苏区的人口数据，还可从区域面积人口数加以推算。闽西地广人稀，1926年人口密度为82人/平方公里，1935年下降到56人/平方公里。\footnote{参见《京粤线福建段经济调查报告书》，铁道部业务调查科编印，1933，第64页；《福建省统计年鉴》，福建省政府编印，1935，第104页。}当然在第五次反“围剿”开始时，这个数据应该会高一些，以此和闽西苏区万余平方公里的实际面积相乘，得数也在百万左右。江西、福建两苏区人口总计仅过300万之数。\footnote{关于中央苏区人口，凌步机著《中央苏区人口数新考》（《中央苏区史研究文集》，中共赣州地委党史办公室编印，1989）有较为详尽的考证，可参阅。}


有限的总人口背后，还有不能不正视的人口比例的失调。苏区建立后，经过连年“围剿”与反“围剿”的征战，青壮年资源损失巨大。当时有材料显示，经过累年的输送军队，1933年初中央苏区乐观的估计“尚有七十万壮丁”。\footnote{亮平：《纪念五一论红军建设中当前的几个重要问题》，《斗争》第10期，1933年5月1日。}为准备第五次反“围剿”，红军又大举增兵，1933年5月，江西扩大红军26520人，兴国一县即达5638人。5月至7月，总计扩大红军约5万人。


第五次反“围剿”开始后，由于出现重大人员消耗，虽然兵源已十分有限，红军仍不得不继续增兵。自1933年8月至1934年7月，中央苏区扩大红军的总人数为112105人。其中1933年8月扩大6290人，9月扩大5868人，10月扩大2214人，11月扩大1958人，12月至1934年1月扩大23258人；1934年2月扩大5865人，3月扩大3344人，4月扩大2970人，5月扩大23035人，6月扩大29688人，7月（至15日止）扩大2457人；其他扩大5467人。\footnote{《一年来扩大红军的统计》，《红星报》第54期，1934年7月22日。}这还不包括各级政权工作人员、工厂工人、前后方夫役等，除去不能参加红军的地主、富农及病残人员，大部分可以参加红军的青壮男子实际都已被征发。据国民党方面占领苏区后的调查，临川、南城、南丰等县壮丁人数均不及全体人口的20\%，南城只占15.49\%，\footnote{《收复区人口数及壮丁比较表》，汪浩：《收复匪区之土地问题》，第49页。}这些只是苏区的边缘区域，中心区域的比例更低。兴国第二区第一、二两保共720人中，壮丁只剩下67人。兴国第二区原有3万人，红军长征后统计只剩1.3万人。\footnote{汪浩：《收复匪区之土地问题》，第50页。}这其中，固有些是逃亡未归者，但大多数是由于军事、政治因素的减员。瑞金、兴国等县仅在后方充任夫役人员死亡数即达百人以上，其中年龄最大者为64岁。\footnote{参见《江西苏区交通运输史》，人民交通出版社，1991，第176～190页相关各表。}


毛泽东1933年底在兴国长冈乡和福建才溪乡的调查，也证实了当时苏区兵源穷尽的现实。兴国长冈乡全部青壮年男子（16岁至45岁）407人，其中出外当红军、做工作的320人，占79\%。才溪乡全部青壮年男子1319人，出外当红军、做工作的1018人，占77\%。\footnote{毛泽东：《才溪乡调查》，《毛泽东农村调查文集》，第351页。}近80\%的比例，已是可征发人员的极限。当然，长冈和才溪是先进乡，不完全反映苏区大多数地区的状况。更具普遍意义的数据显示，1934年年中，中央苏区“红军家属的人口一般的占全人口的一半，在兴国、瑞金、太雷、杨殷、上杭等有些区乡已达三分之二”。\footnote{王首道：《模范红军家属运动》，《斗争》第70期，1934年8月16日。}石城全县人口13.6万余人，其中参加红军者16328人，\footnote{温昌义：《石城红军知多少》，《石城文史资料》第2辑，第55页。}这在苏区尚不能算是很高的比例，宁都全县人口273652人，\footnote{《江西省苏报告》，《红色中华》第41期，1932年11月21日。宁都人口不同调查差别较大，江西省苏1932年5月的调查是204651人。}其中男性138779人，女性133872人，有56304人入伍，参军比例达到全体男性的40.58\%，绝大部分青壮年都参加了红军。大量青壮年男子投入前方，耗去了苏区最宝贵的人力，同时繁重的夫役负担对留在后方的农民和干部也是沉重压力，像福建寿宁县就有“赤卫队一千多人脱离生产”，导致“农村破产（无人耕田种山，粮食恐慌）”。\footnote{《中共福安中心县委关于扩大会议后三个月来的工作报告（1934年1月3日）》，《福建革命历史文件汇集》甲19册，第240页。数字有误，原文如此。}宁都“固村、梅江发生不去当运输的打他的鸡子，竹窄嵊、丸塘乡用捆绑的办法，东山坝，特别是洛口，不去当运输的罚苦工，禁闭；城市、洛口无力担扛的抽担架税，出二毫三毫十毫不等”。\footnote{泽鸿：《宁都参战工作的检阅》，《党的建设》第4期，1932年9月10日。}好在苏区妇女许多都是天足，可以承担相当一部分任务，第五次反“围剿”期间，“运粮食，抬担架基本上是妇女”。\footnote{《赣州军分区干休所座谈记录·苏区的运输和供应》，《回忆苏区交通——苏区交通史料选编》，江西省交通厅史志办公室编印，1987，第139页。}


青壮年男子基本离开土地后，农村劳动力严重缺乏，1934年4月，兴国县的红军家属达61670人，\footnote{《兴国优待红军家属工作》，《红色中华》第203期，1934年6月16日。}留在后方的基本都是老弱病残。虽然苏区中央努力动员更多妇女参加田间劳动，并加紧“调动地主富农举行强迫劳动”，以最大限度地“节省我们工农群众自己的劳动力”，\footnote{《争取决战目前扩红突击的胜利》（社论），《斗争》第60期。1934年5月19日；张闻天：《对于我们的阶级敌人，只有仇恨，没有宽恕》，《红色中华》第193期，1934年5月25日。}但劳动力缺乏仍然成为突出问题。主要因劳动力缺乏，苏区出现大量荒田，长汀水口1932年间“一个区荒田有五千七百担”。\footnote{《团闽粤赣省委三个月工作报告》，《闽粤赣革命历史文件汇集（1932～1933）》，第163页。}1934年初，整个中央苏区荒田达到80万担，“单在公略一县，就有二十八万担”。\footnote{亮平：《把春耕的战斗任务，提到每一劳苦群众的面前》，《斗争》第49期，1934年3月2日。}胜利等县的统计数字也相当惊人，到1934年6月，没有莳禾的田地“胜利有三万四千余担，万太有二万五千余担，洛口吴村区有一千五百余担，博生梅江区二千余担”。\footnote{钟昌桃：《不让一寸土地荒芜》，《红色江西》第4期，1934年6月26日。}


由于后方劳动力缺乏，苏区中央制定的优待红军家属规定难以正常执行，有些地区红军家属由于缺乏劳力，土地被迫弃耕。1934年春耕中，“全苏区大约有二三万担红属的田荒芜着（瑞金全县红五月底止，约有二千七百余担田未莳）……原因除了当地耕田队的消极怠工以外，有些地方确实是因为劳动力缺乏”。\footnote{然之：《把优待红军家属工作彻底改善起来》，《斗争》第66期，1934年6月30日。}万泰县的刘士进，“两个儿子都当红军，分到的六十二担田，前年已荒了十担，今年又荒了六担……现在有好久没有米吃”。\footnote{《铁锤向着窑下区》，《红色中华》第204期，1934年6月19日。}汀州红属由于田地荒芜，1934年夏收前缺粮者达到1575人。\footnote{《汀州市解决红属的困难问题》，《红色中华》第206期，1934年6月23日。}长汀河田区“红军家属有因得不到帮助与优待而做了叫花子”。\footnote{严仲：《三个月扩大红军的总结与教训》，《中央革命根据地史料选编》（中），第664页。}勉强能够执行优红条例的地区，也是不堪重负：“瑞金隘前区，据区土地部的报告，那边每一劳动力每月要帮助红军家属做十六天工。”\footnote{定一：《春耕运动在瑞京》，《斗争》第54期，1934年4月7日。}如此紧张的劳力使帮工时间和质量都难以保证，瑞金红军战士回家后看见家属困境“竟有出眼泪的，同时后方同志看见也就不想去当红军了”；\footnote{《黄沙区的严重现象》，《红色中华》第143期，1934年1月16日。}红军家属的窘境，既影响到前方将士的士气，也使本来就十分困难的扩红工作更形紧张。


人力缺乏对扩红工作影响至巨。由于前线出现大量牺牲，不断补充红军在所必需，但后方人员补给可选择余地愈来愈小，在此背景下，红军动员工作开展艰难：“一般党团员对扩大红军工作是很消沉的，自己也怕当红军。”\footnote{《梁广给全总执行局的报告》，转见刘少奇《反对扩大红军突击运动中的机会主义的动摇》，《斗争》第41期，1934年1月5日。}从1933年下半年开始，每月的扩红指标都难以如期完成，1934年5月，一贯勉力走在前面的江西省只完成计划的20\%。\footnote{《博古同志给李富春同志的信（1934年5月26日）》，《斗争》第63期，1934年6月9日。}这种普遍难以完成计划的情况，虽然当时苏区中央一再以“没有具体的开展反机会主义的斗争”等加以解释，但客观看，人力资源的异常匮乏无疑是主要原因：“上杭县才溪乡，共有二千余人口，在一次一次的扩军突击后，乡里只剩下壮丁七人，还要进行突击，这当然不现实。”\footnote{《李志民回忆录》，第223页。}“太雷县青年男女可以加入少队的只有6135人，而省队部规定它发展到8千人。后来龙岗、石城也发生同样现象。”\footnote{《赤少队突击运动的总结与红五月动员》，《斗争》第57期，1934年4月28日。}这种屡屡出现的发展指标大于实际人数的状况，当然绝不能仅仅用所谓工作失误解释。


一方面是前方需要大量补充，另方面是后方已经罗掘一空，扩红难免不陷入强迫命令的窘境。于都军事部长陈贵公开说：“不用绳捆，有什么办法扩大红军？”\footnote{罗迈：《在粤赣省第一次党代表会议的前面》，《斗争》第34期，1933年11月12日。}对普通群众，有的地方“不去当红军的就封他的房子”。\footnote{潘汉年：《工人师少共国际师的动员总结》，《斗争》第24期，1933年8月29日。}宁化县扩红时，强迫成为主要方式：

\begin{quote}

城市少共市委组织科到群众家去宣传当红军，如发现家里动员对象不在家时，便认为是逃跑了，是“反革命”的了，于是甚至把其家属捉起来。石碧区个别乡召开扩红动员大会时，群众进入会场后，即把门关上，开会动员后让群众“报名”，不肯报名的人，便不准他离开会场。方田区军事部长把不去当红军的群众，派人用梭标解到区苏去。\footnote{《宁化落后的原因在哪里？》，《红色中华》第238期，1934年9月26日。}
\end{quote}


扩红的压力使各级干部日子也不好过，许多干部因完不成任务遭到各种各样的处罚。胜利县“硬要全体干部去当红军，结果‘你不去我也不去’，以至走到‘连干部都上山躲避逃跑甚至个别的自尽的严重现象’”。\footnote{富春：《把扩大红军的突击到群众中去》，《斗争》第38期，1933年12月12日。}江西“龙岗少共县委宣传部长到一个支部去开会，什么都不说，只一声命令，支部书记限三天完成突击，否则杀头，结果这个支部书记找不到出路吊颈死了”。\footnote{《江西省委通讯》第75期，1934年3月5日。}扩红中出现的问题，张闻天当时就有充分的注意：

\begin{quote}

扩大红军工作中，强迫命令是常见的事，如若谁一次开了小差，那就非绑起来不可……这种命令主义，在党的领导者与党员同志中间也是常常发生的，曾经为了要支部同志报名当红军，支部书记将支部同志整晚关在会场上不放的这种奇怪事情……扩大红军中间，如若我们完全采用强迫命令的方法，那必然是群众的登山，群众的反抗。\footnote{张闻天：《关于新的领导方式》，《斗争》第20期，1933年8月5日。}
\end{quote}


可是，在军事紧急而人力又严重不足的状况下，要保证前方的兵员供应，巧妇又如何能奈无米之炊。


普遍存在的强迫扩红大大影响了红军的素质。在扩红过程中，为完成指标，地方经常是“胡乱收罗”，各地扩红“很少注意对质量的选择，只是为着凑人数乱拉，以致扩大新战士，妇女小孩要占五分之一，甚至有地主、富农等阶级异己分子”。\footnote{《中共闽粤赣苏区省委关于紧急动员工作检阅的总结（1932年12月15日）》，《闽粤赣苏区革命历史档案汇编》第1辑，档案出版社，1987，第249页。}于都工人师征召的三千人中，老人、孩子有122人。\footnote{汉年：《工人师少共国际师的动员总结与今后四个月的动员计划》，《斗争》第24期，1933年8月29日。}集中到部队中的人员情况也不乐观：兴国模范师成立时，5161人的成分为：成年占65.8\%，青年占34.2\%，上了岁数的成年人占大多数。\footnote{江西《省委通讯》第5期，1933年6月15日。另外，工人占21.4\%，党团员占38\%，其中党员23.5\%，团员14.4\%，结婚的占72\%，未婚的占28\%。}福建上杭1932年10月扩大红军567人，其中31岁以上的120人，18岁以下的86人，两者相加占总数的40\%。\footnote{《上杭县扩大红军统计表》，《福建革命历史文件汇集》甲19册，第66页。}这一比例随着苏区的不断扩红还在增高。红军总政治部1934年4月统计，红军战士中16岁以下的占1\%，40岁以上的占到4\%，相当部分人员在30～40岁之间。\footnote{李光：《中国新军队》，第279页。}这些人多已有家累负担，体力也不如青年人，以其为主组成的军队，凝聚力和作战力无法不受影响。


由于战斗损失巨大，大量老兵以及指挥员战斗减员，许多新兵来不及必要的培训就不得不投入前线，对部队战斗力也有很大影响。国民党方面资料记载：红军“缺乏训练，且多新兵，常畏缩不前”。有些新兵“仅训练六天，即被解到前方，补充、参加作战”\footnote{《赣粤闽湘鄂北路剿匪军第三路军五次进剿战史》（上），第八章，第59页；《国军五次围剿赣匪崩溃近况》，《军政旬刊》第19、20期合刊，1934年4月30日。}。蒋介石则观察到：

\begin{quote}

现在匪军的精神与从前完全不同，据投诚的土匪说：他们监视军队的政委，也不比从前那样认真了！从前匪内一般政委，的确自己能够上前督率，自己能够身先士卒，不怕死。但是这一般政委，到现在大半都打死了，而新来的一般政委，精神和能力都不够，稍微遇到一点危难的时候，他自己就恐慌的了不得。\footnote{蒋介石：《主动的精义与方法》，《庐山训练集》，第198页。}
\end{quote}


中共方面有关记录也不避讳训练不良的事实，高虎脑战斗中，战斗在前线的红三军团四师新兵众多，“许多人还不会投掷手榴弹”。\footnote{《张震回忆录》（上），解放军出版社，2003，第80页。}时任红五军团第十三师师长的陈伯钧在部队发生逃跑问题后也总结道：“干部太差，我们派去的工作人员不帮助解决实际问题。”红十三师战斗力应属中上，但大量的新兵补充使其进行射击练习时，“成绩非常不好，十人就有九人脱靶”。\footnote{《陈伯钧日记（1933～1937）》，第249、266～267页。}由于缺乏训练，“个别部队在作战中因伤亡（主要是因疾病和掉队）而损失的人数有时竟高达百分之五十，而在老部队中‘正常’减员在百分之十到二十之间”。\footnote{〔德〕奥托·布劳恩：《中国纪事1932～1939》，第50页。}


随着反“围剿”战争的进行，苏区不断被压缩，这使本来就异常缺乏的人力供应雪上加霜，与此同时，为准备即将到来的战略转移，苏区中央又于1934年9月提出一个月内扩充3万红军的目标。\footnote{《关于九月间动员三万新战士上前线的通知（1934年9月1日）》，《红色中华》第229期，1934年9月4日。}面对这样奇高的指标，基层为完成任务，除强迫命令外实在难觅他路，而这更加剧了群众的对立情绪，在红色首都瑞金已有“群众大批逃跑，甚至武装反水去充当团匪，或逃到白区去。瑞金河东区和长胜区都有这样的事情发生”。\footnote{《总动员武装部副部长金维映同志谈扩红动员不能迅速开展的基本原因》，《红色中华》第234期，1934年9月16日。}对于一个在群众中生存的政党和政权，这样的现象出现，真可谓是几多伤心、几多无奈。

\subsection{物质资源}


中央苏区大部分位于山区，面积在4万平方公里左右，经济发展水平较低，近代工业很少，生活必需的工业品多以苏区出产农产品换取。第五次“围剿”期间，国民党实行严密封锁政策后，苏区物质资源的不足日渐显露，许多当年曾在苏区者都能回忆起苏区物质资源的困窘局面：“红军供应困难，吃没吃的，穿没穿的，打仗缺乏弹药，加上长期转战，部队得不到休整，个个都像叫花子一样”。\footnote{《王平回忆录》，解放军出版社，1992，第52页。}这其中，莫文骅的两段回忆颇为生动：

\begin{quote}

有一天，我师一位团长对他的小勤务员发火，勤务员不高兴地说：“你再发火，我明天不给你手巾洗脸！”原来，这位团长已经很久没洗脸巾了，常借用勤务员的洗脸巾。

一次，我师有一位连长在防御前沿阵地趁黑夜偷越警戒线，企图逃跑，哨兵发现后，立即令他站住，他还算听话，没有继续越界……保卫干部把那位连长押到师部，我亲自进行审问：“你是起义过来当红军的，由士兵升为连长，组织很信任你，现在生活虽然艰苦，但革命是有前途的，你为什么要逃跑？”他没有马上回答我，只是哭，最后难过他说：“我并不怕死，也不怕苦，但是……”他说到这里停了一会儿，才慢慢说出事情的经过：“前天早晨，我召集全连战士训话时，有几个战士指着我的裤子笑，我低头一看，原来自己的裤裆破了，裂开了一个大口子，露出了下身，感到万分羞耻，急忙解散队伍。晚上，睡不着，认为自己不好再继续担任连长了，便决定逃跑。现在冷静一想，真是罪过，后悔莫及！”对这位连长的问题，本应按军法严处，考虑他为革命立过功，且态度较老实，认罪较好，故免除他死刑，判5年徒刑，发给一条旧军裤，送他到后方服刑。\footnote{《莫文骅回忆录》，解放军出版社，1996，第231～232页。}
\end{quote}


苏区物资的困乏局面，实行封锁的国民党方面自然最乐意看到。早在第五次“围剿”之初，他们就报告：“匪军缺乏食盐棉衣等物，逃亡甚多。”\footnote{《吴奇伟电蒋中正陈诚棠荫一役共军伤亡甚巨等匪情侦报》，蒋中正文物档案002090300066401。}随着“围剿”的进行，其封锁效果日益显现，国民党军报告中明显可看到幸灾乐祸的成分：“我方封锁匪区以来，匪方对于生活和给养上，已感受重大的困难，现于俘匪的鸠头鹄面，衣不蔽体，这已是一幕活写真了。”\footnote{《国军五次围剿赣匪崩溃近况》，《军政旬刊》第19、20期合刊，1934年4月30日。}


就维持生存最基本的需要而言，粮食和油盐是不可或缺的必需品。中央苏区所在的赣南属于粮食略有出超区域，而闽西地区粮食则为入超。不过，赣南的出超有限，且建立在当地农民节衣缩食——“乡间多用杂粮佐餐”\footnote{《兴国粮食问题应如何解决以期军需民食得兼筹并顾请公决案》，江西省兴国县档案馆藏档131/2-5-2/58。}上，基础不是十分牢固。清人曾写道：“赣亡他产，颇饶稻谷，自豫章吴会，咸仰给焉。两关转谷之舟，日络绎不绝，即俭岁亦橹声相闻。盖齐民不善治生，所恃赡一切费者，终岁之入耳，故日食之余，则尽以出粜，鲜有盖藏者。”\footnote{天启：《赣州府志》卷三《舆地志·土产》。}福建是缺粮省，缺粮差数达30.3\%，闽西也属缺粮地区，时人有言：“闽西、闽东又多为贫瘠之区，米谷产量甚少，人民素赖甘薯为主要原料。”\footnote{陈颖光：《福建粮食统制之研究》，《闽政公余非常时期合刊》第2期，1937年9月20日。}不过，更加具体的材料则显示，闽西缺粮尚不十分严重，正常年份下，勉强维持平衡尚有可能，下表显示的是1935年闽西苏区主要地区所在的第七区稻谷产需状况：

\begin{figure}[htbp] % htbp 控制图片排版位置优先顺序
    \centering
    \includegraphics[width=0.6\textwidth]{6.jpg}
    \caption{福建省第七区二十四年各县收获稻谷数目及盈亏概况表}
    \label{fig: 福建省第七区二十四年各县收获稻谷数目及盈亏概况表}
\end{figure}

从图\ref{fig: 福建省第七区二十四年各县收获稻谷数目及盈亏概况表}看，以盈亏相抵，第七区粮食总体尚略有盈余。当然，1935年是粮食生产偏好年景，该区人口又大幅减少，不一定有足够的代表性。正常状况下，闽西地区粮食略显不足的判断应可成立。


赣南、闽西成为统一苏区区域后，这里的粮食供应从原来并不宽裕的状况转为紧张。尤其军队急剧增加，大批人员涌入，战争及土地关系的不稳定又影响着生产的发展，使得粮食问题愈益凸显。


祸不单行，1932、1933年中央苏区恰逢灾年，加上政治变动及战争的影响，粮食收成很不理想。1932年中央苏区各县粮产量普遍只有正常年份的60\%左右（最高的新泉、长胜为76\%和75\%，最低的兴国、瑞金、万泰、博生都只有50\%），\footnote{定一：《两个政权——两个收成》，《斗争》第72期，1934年9月23日。}1933年虽有增加，但仍未恢复到苏维埃革命前的水平。苏区建立后不久，就出现缺粮问题，部分缺粮严重地区如富田春荒时“甚至吃草”。\footnote{《团江西省委关于目前紧急任务的决议》，《江西革命历史文件汇集（1932年）》（一），第63页。}到1933年春，中央苏区开始严重缺粮，陈云在1934年间谈道：“去年青黄不接时，因为某些地方缺乏粮食与缺乏全盘及时的调剂，再加上奸商富农的抬价及囤积，曾经威胁了我们。”\footnote{陈云：《为收集粮食而斗争》，《斗争》第45期，1934年2月2日。}当时，“许多地方发生粮食恐慌（如博生城市与赣县、兴国个别区乡），甚至有因为粮食找不到出路的。在博生县有个群众吃药自尽。在兴国黄塘区因为两升米的问题，一个男子把他的兄嫂活埋了”。\footnote{《中共江西省委红五月工作总报告（1933年6月17日）》，《江西革命历史文件汇集（1933～1934年）》，第153页。}应该指出的是，这种粮食恐慌不仅出现在苏区辖境，国民党控制的赣南区域也有反映，1933年5月，陈济棠曾向蒋介石转呈余汉谋的电报：内称“赣南米荒，因无食而吊颈投河者日有所闻。至金坎、快顺、崇义一带均食泥饼，此泥饼亦售六百一斤，职尝之为香粉饼。现虽已派彭李两团到万安掩护谷船来赣，但以水浅每日只行十里，须月余才能到达”。11月，又报告：“赣南今秋水旱失时，晚稻收获仅及往岁十分之三。”\footnote{《陈济棠电蒋中正赣南米荒情形（1933年5月25日）》，蒋中正文物档案002080200089151；《陈济棠电蒋中正转电余汉谋所述赣南米荒情形（1933年11月1日）》，蒋中正文物档案002080200130022。}可见这一时期灾害确实是影响赣南地区粮食供应的主要因素。


随着第五次“围剿”开始，苏区区域压缩，国民党封锁日益严密，粮食问题更形严重。1934年初已出现“红军部队及政府机关食米不够供给”\footnote{《第二次全苏代表大会主席团中国共产党中央委员会关于完成推销公债征收土地税收集粮食保障红军给养的突击运动的决定》，《中共中央文件选集》第10册，中共中央党校出版社，1991，第82页。}的严峻局面。2月召开的中央粮食会议透露了形势的严重：“粮食问题已经非常严重的摆在我们面前，谷价到处高涨，有些地方如会昌、瑞金、博生、于都等处已涨到七八元一担。应该收集的土地税和公债谷子还差着很巨大的数目。”\footnote{《中央粮食会议纪要》，《红色中华》第146期，1934年2月6日。}3月，赣县米价更涨到“十七块多一担”。\footnote{《各地米市简报》，《红色中华》第158期，1934年3月6日。}为节省物资，中央苏区政府号召广泛开展节省运动，规定：“各级政府、红军后方机关、国家企业、学校等每人每天减发食米二两”。随后，又提出苏区每人每月“节省三升甚至三升以上的米来供给红军”。\footnote{《人民委员会批准减少食米的请求》，亮平：《红军等着我们的粮食吃》，《红色中华》第168期，1934年3月29日；第200期，1934年6月9日。}由于供给不足，实际能够发放的粮食数量远远低于定量：“党政机关……每人每天只10小两（1斤为16两）粮食，分成两顿吃。”最艰难的时候，红军战斗部队甚至“每天只能吃八两至十两”。\footnote{《来自井冈山下——罗通回忆录》，东方出版社，1996，第102页；《杨成武回忆录》，解放军出版社，1987，第23页。}与此同时，国民党军又趁火打劫，专门制定规章，组织割禾队，怂恿区外农民“由驻军率领，冲入匪区”抢割稻禾，要求不能收割时，“亦应设法予以销毁”。\footnote{《国民政府军事委员会委员长南昌行营处理剿匪省份政治工作报告》（下），第11章，第48页；蒋介石：《电剿匪区内各军政长官为掩护民众进入匪区时应尽量收割否则应予销毁》，《军政旬刊》第35期，1934年9月30日。}福建宁化1934年6月报告，由于国民党军的破坏和抢夺，“全县到现在止共损失边区谷子一千七百余担”。\footnote{《宁化边区损失许多米谷》，《红色中华》第202期，1934年6月14日。}永丰石马被国民党军收去谷子达2076担。\footnote{《孙连仲致蒋介石电（1934年9月9日）》，《军政旬刊》第35期，1934年9月30日。}


为应付严重的粮食问题，1933年12月，苏维埃中央决定成立粮食人民部，专门处理粮食问题，收集粮食被提到“国内战争中一个残酷的阶级斗争”\footnote{陈云：《为收集粮食而斗争》，《斗争》第45期，1934年2月2日。}的高度。虽然采取了多种措施，粮食收集仍遇到许多困难。1934年2月的粮食突击运动中，“于都预定计划收四万四千担谷子，现在收到百分之十，胜利收谷也只有百分之十三”。\footnote{毛泽覃：《为全部完成粮食突击计划而斗争》，《斗争》第49期，1934年3月2日。}于都古田区、黎村区甚至“没有收到一粒谷子”，该县粮食部长明确认为：“于都群众没有谷子。”\footnote{《于都突击运动中的严重问题》，《红色中华》第155期，1934年2月27日。}在收集粮食遇到困难时，一些地区为完成任务强行摊派，造成政府与群众间关系紧张，胜利和于都等地发生群众“要捉突击队员”并“向区苏请愿”\footnote{毛泽覃：《为全部完成粮食突击计划而斗争》，《斗争》第49期，1934年3月2日。}的严重事件。


粮食出现困难同时，兵员仍在继续增加，不断增加的兵员要求更多的军粮供应。为此，1933至1934年间，中央苏区不得不先后向群众借谷24万担和60万担。应该说，苏维埃政权一直在尽可能不过多加重农民负担，1932年，赣西北的江西万载等地苏维埃政权曾向群众借谷，湘鄂赣省苏维埃政府发现后严加制止，强调：“今后无论何种用费，困难到任何程度，都不得向群众借一斗米或一串或几百钱，尤其是‘预征税收’的国民党工作方式，更不容许应用。”\footnote{《湘鄂赣省苏维埃政府通令（1932年6月19日）》，《湘鄂赣革命根据地文献资料》第2辑，第269页。}此时，中共的借谷决定确实是不得已而为之。由于粮食供应本已相当困难，借谷任务除兴国等少数地区外，大部分地区完成得相当艰难。苏区中央严令各地“派出的突击队不到任务完成不能调回”。\footnote{《中央组织局、人民委员会关于粮食动员的紧急指示》，《红色中华》第209期，1934年6月30日。}《红色中华》发表文章强调：“目前我们是处在决死的战斗中……如果没有二十四万担粮食，我们的红军不能作战。不能吃饱肚子，就不能维持生活。”\footnote{定一：《动员二十四万担粮食是目前我们第一等的任务》，《红色中华》第210期，1934年7月5日。}在战争条件下，政府对苏区粮食资源的挖掘可以说已经到了最大限度。


和粮食的严重局面比，一些依赖输入的日用品更形紧张，苏区的食盐、布匹、煤油、药材等生活必需品极端匮乏，据毛泽东1933年底调查：

\begin{quote}

暴动前平均每人每两年才能做一套衫裤，暴动后平均每人每年能做一套半，增加了百分之二百。今年情形又改变，因为封锁，布贵，平均每人只能做半套，恢复到暴动前。暴动前一套单衣服值十八毛（十五毛布，三毛工），去年每套二十一毛（十七毛布，四毛工），合大洋一元半，今年每套三十四毛（三十毛布，四毛工），合大洋二元四角。反革命使我们的衣服贵到如此程度！……暴动前每人平均每月吃盐一斤，今年十一月每人每月只吃三两二钱，即暴动前五个人的家庭月吃盐五斤者，今年十一月只吃一斤。\footnote{毛泽东：《长冈乡调查》，《毛泽东农村调查文集》，第350页。}
\end{quote}


这还是第五次“围剿”初期的情况，随着战争的进行和封锁的加紧，问题愈加严重。尤其是作为生存必需品的食盐。赣南食盐向来依赖外来的淮盐、潮盐、广盐补给，多“由土民担谷往赣北或广东交界处交换而来”。\footnote{《陈诚罗卓英电蒋中正据夏楚中据投诚匪兵供称伪第三军团师辖有三团所携食盐已尽》，蒋中正文物档案002090300074117。}据1930年《申报》报道：“平均潮属每年所产之盐，销入闽赣外省者占十分之七。”\footnote{《粤盐推销闽赣之障碍》，1930年7月21日《申报》。}闽赣对外来食盐的依赖，可见一斑。由于国民党军实施封锁，严禁食盐输入，导致苏区食盐价格飞涨，1933年底，一元大洋只能买盐一斤多，相当于红军到来前的1/4。\footnote{毛泽东：《长冈乡调查》，《毛泽东农村调查文集》，第308页。}很多家庭买不起盐，只能淡食。1934年初，食盐价格继续上涨，私人商店每元只能买盐半斤多。中央粮食调剂总局出售的平价食盐为每元一斤，但仅在瑞金设点销售，且“不能供给非红军家属的需要”。\footnote{《粮食调剂总局举行廉价》，《红色中华》第158期，1934年3月6日。}1934年年中，随着前方战事不利，物价再现暴涨：“米每元五斤，盐每元一两五钱，柴每斤二角。”\footnote{《汤恩伯1934年6月2日致蒋鼎文电》，《中华民国史档案资料汇编》第5辑第1编军事（4），第204页。}为解燃眉之急，苏区不得不“大举进行熬硝盐的事业”。\footnote{《争取决战目前扩红突击的胜利》（社论），《斗争》第60期，1934年5月19日。}迫不得已时，“把厕所底下的土，挖出来熬盐，甚至用死人墓下的土熬盐”。\footnote{王观澜：《中央苏区的土地斗争和经济情况》，《回忆中央苏区》，江西人民出版社，1981，第320页。}硝盐质地不纯，成分也有别于食用盐，食用硝盐后，中毒事件屡见不鲜。


物资奇缺，导致价格飞涨，财政金融部门为弥补财政赤字，大量发行纸币，苏区中央有关文件承认：“在长期国内战争的条件之下，增发纸币常常是弥补财政收支不敷的一个办法。”\footnote{《关于苏维埃经济建设的决议——第二次全国苏维埃代表大会通过（1934年1月）》，《革命根据地史料选编》（上），第169页。}在此背景下，严重的通货膨胀难以避免。瑞金城区有一商人“以三元苏维埃纸票收买现洋一元，后来又将所收的现洋以三元五角价钱出卖给别人”，此人被定性为进行“捣乱经济、操纵金融的反革命活动”，以反革命罪枪决。\footnote{《国家政治保卫局又逮捕了四个反革命》，《红色中华》第152期，1934年2月22日。}其实，这种倒买倒卖现象是苏区金融秩序失控的真实反映。


由于第五次“围剿”以来国民党军采取稳扎稳打战术，各部不轻易跃进，红军擅长的诱敌深入打歼灭战战术难以发挥，无法缴获到大批武器弹药，这使红军主要的武器来源受到限制。红军兵工厂本身的制造能力、技术能力、管理能力与材料储存又十分有限：“不能制造枪、弹……枪枝只能修理，子弹只能翻造。”\footnote{马文：《忆中央苏区红军兵工厂》，中国兵器工业历史资料编审委员会编《土地革命战争时期军工史料》，中国兵器工业总公司，1994年内部出版，第112页。}翻造的子弹，即将打过的弹壳装上新造弹头而成，这种子弹因为“装的是土造的硝盐，是从厕所墙壁上刮下来的尿碱熬成的，燃烧速度慢，动力不足。弹丸是用电线拧成的一坨铁蛋蛋，不能啮合膛线，初速很低，所以打出去之后在空中折跟斗”。\footnote{《耿飚回忆录》（上），江苏人民出版社，1998，第146页。}刘少奇当时曾写道：“兵工厂做的子弹，有三万多发是打不响的，枪修好了许多拿到前方不能打，或者一打就坏了。”\footnote{刘少奇：《论国家工厂的管理》，《斗争》第53期，1934年3月31日。}子弹如此，手榴弹和刺刀也差不多：“手榴弹都是我兵工厂造的，质量太差，落地只炸成两三块，杀伤力很低。敌人上来后，只有拼刺刀了，可有的刺刀也不顶用，捅几下便弯了。”\footnote{《张震回忆录》（上），第80页。}


随着战争的持久进行，红军作战物资消耗严重，枪械、弹药供给越来越困难。国民党方面在1934年年中侦察到：“伪一师每连仅有士兵三四十名，子弹每枪约六七排，都是土造，连续射击不得超过五发，以上则炸裂。”福建方面一些地方红军枪弹缺乏，“多持标枪扁担”。\footnote{《汤恩伯1934年5月11日致蒋鼎文电》、《卢兴邦1934年8月13日致蒋鼎文电》，《中华民国史档案资料汇编》第5辑第1编军事（4），第190、239页。}据统计，1934年8月，红一军团共有步马枪7268支，三军团5385支，五军团4000支左右，九军团1830支，\footnote{《朱德关于补足军械器材集中物资致各部队电（1934年9月9日至20日）》，中国人民解放军历史资料丛书《后勤工作·文献（1）》，解放军出版社，1997，第403～405页。}几个主力军团枪支数不足两万。红军一个主力连的武器配置为：步、马枪共41支，步、马枪弹2025发，驳壳枪2把，驳壳枪弹9发，轻机枪1挺，轻机枪弹60发，手榴弹30颗，\footnote{《红军人马武器弹药月终总报告表（1934年8月18日）》，石叟档案1008·564/2138，中国社会科学院近代史研究所藏。}弹药筹备几乎难以支撑一个大的战役。重武器更是无法和国民党军相比，当时，“敌人有一千五百门迫击炮和野战炮，而我们只有二十几门，就是这些炮又由于缺乏弹药多数已无法使用。机关枪的比例虽然不那么悬殊，但也至少是二十、三十比一”。\footnote{〔德〕奥托·布劳恩：《中国纪事1932～1939》，第56页。}正是由于吃准了红军作战资源缺乏、无力攻坚的弱点，蒋介石可以放心大胆地使用碉堡战术，如他自己所解释的：“倘使他们有炮，又有炮弹，那么碉堡自然是没用的。”\footnote{《蒋委员长出巡纪要》，《军政旬刊》第37、38期合刊，1934年10月31日。}

\section{经济力的挖掘}

\subsection{私人经济发展的限制}


作为在战争中发展起来的红色区域，中央苏区资源匮乏，经济力也不容乐观。美国记者斯诺观察长征以后中共在西北地区建立的苏区后写道：

\begin{quote}

不论中国共产主义运动在南方的情况如何，就我在西北所看到的而论，如果称之为农村平均主义，较之马克思作为自己的模范产儿而认为合适的任何名称，也许更加确切一些。这在经济上尤其显著。在有组织的苏区的社会、政治、文化生活中，虽然有一种马克思主义的简单指导，但是物质条件的局限性到处是显而易见的。\footnote{〔美〕埃德加·斯诺：《西行漫记》，董乐山译，三联书店，1979，第193页。}
\end{quote}


较少具有政治偏见的斯诺在这里依然表现出了他作为优秀记者的敏锐和洞见，确实，如其所言，在贫窭的苏区经济条件下，中共秉持的理念及现实环境的需要，都使他们将平均作为其推行政策的极其重要的逻辑基础，而中共所迅速获得的支持，和这一理念密不可分，尽管平均也是一把双刃剑，在平均、平等、效率、发展等问题之间，总会纠缠着一些难解之结，但中共在力所能及的范围内，还是尽可能地在苏区建立起了一套较有效率的社会经济体系，将社会财富的挖掘、调动、分配效用发挥到了极致，以有限的资源，支撑起长达数年的持续战争。


苏区原有的经济基础本就十分薄弱，又不断处于战争破坏之中，经济的维持可谓举步维艰。根据毛泽东在二苏大的报告，中央苏区经济建设的原则是：“进行一切可能的与必须的经济方面的建设，集中经济力量供给战争，同时极力改良民众的生活，巩固工农在经济方面的联合，保证无产阶级对于农民的领导，造成将来发展到社会主义的前提与优势”。\footnote{毛泽东：《中华苏维埃共和国中央执行委员会与人民委员会对第二次全国苏维埃代表大会的报告》，《苏维埃中国》，第278页。}战争、生活是这一原则的关键词，同时，为社会主义准备发展条件则决定着苏维埃经济的性质，尽管对苏区而言，这还比较遥远。


对于中国苏维埃经济所属阶段的认识，中共和共产国际都经历了一个曲折的过程。初期的土地共有、集中农场及对商人和私人资本的打击，反映了在经济上一举过渡到社会主义的愿望。随着这些政策恶果的显现，共产国际也在不断调整其对中国苏维埃经济社会政策的指导，1933年3月，共产国际执行委员会致电中共中央，指示：

\begin{quote}

你们应高度重视苏区的经济政策问题。认真研究和制定相应的指示。请考虑我们提出的以下设想。恶化的经济状况迫切要求进一步明确和修正我们在发展和鼓励生产、活跃市场关系和商品流通方面的经济政策和专门措施。无重大原因，要避免重新分田，特别是在老区。分田之后农民应当拥有固定的土地，只没收那些参与反革命活动的富农的生产资料。要针对雇农制定临时规定，并对中农和富农的雇佣条件加以区分。以往给富农提供份地的条件依然有效，不能把这理解为禁止富农出租和买卖在苏维埃管制下的土地和雇佣劳动力……要鼓励手工业生产和贸易活动。这应该在税收政策和为手工业企业徒工和工人制定的临时规定上有所反映。同时要记住，社会立法和工会工作等都应符合经济实力和红军斗争的利益。\footnote{《共产国际执行委员会政治书记处给中共中央的电报（1933年3月19～22日）》，《共产国际、联共（布）与中国革命档案资料丛书》第13册，第354～355页。}
\end{quote}


1933年底，王明撰文批评了过分与富农斗争的行为，并指出：“有一种现象阻碍苏区农业生产的振兴，这就是再三不断地重新分配土地。”“有些同志抱定‘左’的立场，其表现就是他们‘害怕’资本主义。”“必须明白，我们暂时还不能在中国苏区内消灭资本主义，而只是准备将来消灭资本主义的一切前提和条件。”“目前苏维埃并不提出自己力量还不能完成的任务，就是说还不从苏区经济中铲除资本主义，而是利用它（在苏维埃政权机关所能做到的范围内）以谋振兴苏区的经济生活。”文章批评苏区劳动法“左”的错误，指出其“对于富农、中农、中等的和小的手工业主惯用工人的条件，没有分别看待。这一点在某些苏区的苏维埃的工作实际中，已经产生了不良后果”。“要使苏区商品流通兴旺起来，首先就要坚决实行贸易自由的原则，同时要反对投机买卖及尽力保证红军经常供给。共产党和苏维埃政府应当禁止经济政策方面的一切过火办法，这些办法事实上只不过使苏区经济状况变坏。”\footnote{王明：《中国苏维埃政权底经济政策（1933年11月）》，《王明选集》第3卷，日本汲古书院，1973，第30～39页。}


随着苏区的巩固发展，苏维埃的经济政策确实面临着一些不得不认真考虑的问题。打土豪财政的结束，相对正规的财税制度的建立，使发展生产成为保障整个国家机器运转的必要条件，对此中共领导人不能不认真应对。1933年5、6月份，张闻天、博古相继发表文章，根据莫斯科的指示对经济政策实施调整，一方面继续强调：“我们党的任务是在集中苏区的一切经济力量，帮助革命战争，争取革命战争的胜利，在这中间巩固工农在经济上的联合，在经济上保证无产阶级的领导，造成非资本主义（即社会主义）发展的前提和优势。”同时鉴于苏区的实际状况，客观指出：“苏维埃政府现在还是非常贫困，它没有足够的资本来经营大规模的生产。在目前，它还不能不利用私人资本来发展苏维埃的经济。它甚至应该采取种种办法，去鼓动私人资本家的投资。”“苏维埃政权同某些资本家可以订立协定，甚至给他们以特别的权利，使他们发展他们的企业，扩大他们的生产，这对于苏维埃政权的巩固是有利的。”\footnote{洛甫：《论苏维埃经济发展的前途》、《苏维埃政权下的阶级斗争》，《斗争》第11期，1933年5月10日；第14期，1933年6月5日。}文章强调对地主和资产阶级要区别对待，“苏维埃政权对于这两个阶级的态度不应该一视同仁的，对于地主我们采取坚决无情的消灭地主这一阶级的政策……至于对于资产阶级，那就完全是另外一种情形……消灭资产阶级的政策是不适当与不能采用的，所以，‘左’的空谈消灭资产阶级是有害的”。\footnote{博古：《论目前阶段上苏维埃政权的经济政策》，《斗争》第16期，1933年6月25日。}


应该说，从原则上看，苏区对私营工商业一直采取了保护和奖励政策。1929年3月，红四军在汀州发布《告商人及知识分子》，明确宣布：“共产党对城市的政策是：取消苛税杂税，保护商人贸易。在革命时候对商人酌量筹款供给军需，但不准派到小商人身上……普通商人及一般小资产阶级的财物，一概不没收。”1930年3月，闽西苏维埃政府颁布《商人条例》规定：“商人遵照政府决议案及一切法令，照章缴纳所得税者，政府予以保护，不准任何人侵害。”“所有武装团体，不得借口逮捕犯人，骚扰商店。”\footnote{《闽西苏维埃政府布告第九号》，《福建革命历史文件汇集》甲14册，第113页。}1931年8月，闽西苏维埃政府针对杭武第六区发生擅将商人货物拘押、拍卖事件发布《允许商人自由贸易问题》的通知，严厉告诫不要实行自我封锁的“自杀政策”，指出：

\begin{quote}

商人虽是剥削分子，但在这土地革命时期中，是以铲除封建剥削为主要任务，对于商人，当然不能够照对豪绅地主一样加以打击，所以无论何种商人，只要他在不违反苏维埃法令（如劳动法，土地法等），不勾结敌人作反革命，不操纵与垄断经济，苏维埃政府是允许商人在苏区内营业自由的。
\end{quote}


通知强调这种随意拘捕商人、没收商品的行为，“不但给敌人以造谣破坏的机会，而且各地商人势必都不敢到苏区内营业，油盐布匹都没有买，是不待敌人来封锁我们，而我们先自己封锁自己，这就是叫做‘自杀政策’”。通知要求各地认识到：“经济政策中的‘自由贸易’是目前一个很关重要、很要注意的问题，无论赤色白色区商人我们绝对不要随便去打击他！”\footnote{《闽西苏维埃政府通知第93号》，《福建革命历史文件汇集》甲15册，第175～176页。}这一通知一方面反映了苏维埃政权对商业贸易的维护态度，另一方面也显露出以社会主义为发展方向的苏维埃各方在执行这一政策时的抵触。事实上，由于商人被认为是剥削分子乃至敌对阶层，对于基层政权而言，他们很难分清楚其与地主间的区别，而筹集款项和物质资源的需要也很容易使打击商人成为现实。所以，在苏维埃的商人理论和实际执行之间有一个很大的落差，苏区商业一度陷入极度萧条的境地。


中共大批干部进入中央苏区后，循着正规化、长期化的思路，对维护苏区经济的活力给予较多重视。1932年1月，苏维埃临时中央政府颁布《工商业投资暂行条例》，“以鼓励私人资本的投资”。条例规定：“凡遵（守）苏维埃一切法令，实行劳动法，并依照苏维埃政府所颁布之税则，完纳国税的条件下，得允许私人资本在中华苏维埃共和国境内自由投资经营工商业。”“无论国家的企业、矿山、森林等和私人的产业，均可投资经营或承租承办，但须由双方协商订立租借合同，向当地苏维埃政府登记，但苏维埃政府对于所订合同，认为与政府所颁布的法令和条件相违反时，有修改和停止该合同之权。”\footnote{《中华苏维埃共和国临时中央政府关于工商业投资暂行条例的决议》，《中央革命根据地工商税收史料选编》，福建人民出版社，1985，第70页。}为发展苏区的社会经济，充裕苏区的经济实力，苏维埃政府除鼓励私人投资外，还于1932年8月颁布《矿产开采权出租办法》、《店房没收和租借条例》，规定私人资本可以向苏维埃政府承租矿产开采、店房、作坊等。是年9月13日，中央政府财政部发出训令，要求各地“必须严厉执行经济政策……注意检查各地政府有无破坏经济政策的行为，如胡乱没收商店，乱打土豪，限制市价，随便禁止出口等，如发现有这些行为，必须予以严厉纠正或处分”。\footnote{《目前各级财政部的中心工作》，《中央革命根据地工商税收史料选编》，第128页。}


对于调整后的私人经济的政策，1934年1月，毛泽东在二苏大明确谈道：

\begin{quote}

苏维埃对于私人经济，只要不出于苏维埃法律范围之外，不但不加阻止，而且是提倡的奖励的。因为目前私人经济的发展，是苏维埃利益的需要。关于私人经济，不待说现时是绝对的优势，并且在相当的长期间也必然还是他的优势。

尽可能的发展国家企业与大规模的发展合作社，应该是与奖励私人经济发展同时并进的。\footnote{毛泽东：《对第二次全国苏维埃代表大会工作报告》，《苏维埃中国》，第281页。}
\end{quote}


中共中央的这些规定在一定程度上产生了效果，如在中央苏区周边的中心城市赣州，“广裕兴老板曾伟仁在赣州开设广裕兴（经营百货）以及明为剧院暗做生意的新光剧院等，都有广东军阀入赣驻军李振球（系广东一军副军长兼一师师长，是同乡关系）做后台”。中共“抓住军阀的贪官发财及其内部矛盾，通过在广裕兴商号工作的同学刘××（外称刘少老板，是广东人）的关系做通老板曾伟仁的工作，在默林支部书记丁友生家，进行秘密商谈交换货物，经过协商双方同意互派代表，并就物资的交换等事宜达成了协议。广裕兴派出代表钟先请长驻江口（历时半年左右）收购钨砂和稻谷。江口分局派出刘东门生，在赣州以开设小杂店为名，长驻赣州，为苏区秘密联络和采购物资”。\footnote{黄谋祝：《苏区时期赣县反经济封锁斗争》，《赣县文史资料》第1辑，政协赣县文史资料委员会编印，第60页。}不过，由于大规模战争接连到来和国民党封锁的日益紧密，客观条件限制使其难以收到吸引投资和维护私营工商业的功效，而前期对工商业打击的负作用及各级干部对私营工商业根深蒂固的排斥态度，也使苏区内的私人投资仍然呈现日益萎缩的趋势。

\subsection{合作社事业的发展}


私人商业式微，代之而起的是苏维埃政权投入心力的合作社，合作社事业的发展是中央苏区经济社会变化中最值得注目的现象之一。合作社是由群众出资出物入股经营，带有较多互助成分的集体性质事业，被认为是在生产、供给、分配战线上“战斗地团结工农，动员群众的经济组织”。\footnote{亮平：《目前苏维埃合作运动的状况和我们的任务》，《斗争》第56期，1934年4月21日。}按照张闻天的说法：“在苏区内生产与消费的合作社，不是资本主义的企业，因为资本家与富农的加入合作社是完全禁止的。这是一种小生产者的集体的经济，这种小生产者的集体经济目前也不是社会主义的经济，但是它的发展趋向将随着中国工农民主专政的走向社会主义而成为社会主义的经济。”\footnote{张闻天：《论苏维埃经济发展的前途》，《斗争》第11期，1933年5月10日。}


细究起来，苏区合作社发展风潮和当时全国范围内风起云涌的合作社运动有相通之处，而苏区内合作社表现出的强烈的政治性则为其他地方不易见到。1930年代前后，包括国、共两党在内的多种社会政治力量都在进行合作社运动的尝试，比较而言，中央苏区的合作社通过政治力量组织、推动，其发展速度和效率都相对较高。在一些发展较好的合作社中，有比较严密的监督、运作制度。瑞金“武阳石水乡合作社实行严格的监督制度，五天查一次帐，社中买猪给社员饲养，社员替社中工作人耕田，使社员对合作社取得密切的连系”。\footnote{《中央苏区消费合作社第一次代表大会纪盛》，《红色中华》第133期，1933年12月8日。}福建上杭才溪区消费合作社“卖货由社员大会决定”，“能够发动社员来经常监督工作，审查委员会经常开会，所以没有发生贪污的事情”。\footnote{《一个模范的消费合作社》，《红色中华》第139期，1934年1月1日。}


中央苏区的合作社最早于1928年10月出现在赣西南的东固，后在赣南、闽西推广，主要为流通领域的消费合作社。1932年4月，临时中央政府颁发《合作社暂行组织条例》，规定“合作社组织为发展苏维埃经济的一个主要方式”，“由工农劳动群众集资所组织”。\footnote{《中华苏维埃临时中央政府关于合作社暂行组织条例的决议》，《中央革命根据地史料选编》（下），第573页。}1933年，临时中央政府国民经济部成立合作社指导委员会，专门负责指导合作社的发起、建设、监督管理。到1933年8月，共有消费合作社417个，社员82940人；粮食合作社457个，社员102182人；生产合作社76个，社员9276人。股金总额216095元。\footnote{亮平：《目前苏维埃合作运动的状况和我们的任务》，《斗争》第56期，1934年4月21日。}


1933年8月先后召开的中央苏区南部十七县和北部十一县经济建设大会进一步确定了发展合作社事业的方针，此后，合作社在中央苏区获得迅速发展。瑞金在8月以前，只有消费合作社社员9000人，股金1.1万元，9月份，社员就增加了5300人，股金增加了5500元；粮食合作社在8月前很少，9月份社员增加到6800人，股金增加到1800多元。兴国在大会以后的一个月中，消费合作社社员增加14600人，粮食合作社社员增加15000人。\footnote{《中华苏维埃临时中央政府关于合作社暂行组织条例的决议（1932年4月12日）》，《革命根据地经济史料选编》（上），第87页。}整个中央苏区到1934年2月，消费、粮食、生产三种合作社已“发展到二千三百余的社数，五十七万余的社员和六十余万元的股金”，\footnote{亮平：《目前苏维埃合作运动的状况和我们的任务》，《斗争》第56期，1934年4月21日。}各方面指标都有了近三倍的增长。消费合作社包括17个县总社，两个省总社（江西、福建）及中央苏区的消费合作总社（1933年12月17日正式成立）。


苏区时期，最成功的合作形式当推农业互助的各种组织。苏维埃革命后，劳动力和耕牛缺乏成为新分得土地农民的一个普遍问题，为尽可能利用有限资源，中央苏区总结群众换工互助的传统经验，倡导组织耕田队、劳动互助社和犁牛合作社，以解决人力、畜力缺乏问题。


耕田队最早出现在上杭才溪区，是农村劳动互助的初级形式，苏区农民在个体经济的基础上，为调剂劳动力、恢复农业生产，以自愿、互利原则创办。耕田队的经验得到苏维埃政权高度重视，并在此基础上，指导发展出劳动互助社。劳动互助社在不变更个体所有制的前提下，调整生产关系，一定程度上打破以一家一户为生产单位的界限，对于战争环境劳动力奇缺的苏区恢复和发展农业生产起了重要作用。1933年底，兴国长冈乡的劳动互助社，“四村每村一个，除红属外凡有劳动力的，十分之八都加入了。全乡社员三百多”。\footnote{毛泽东：《长冈乡调查》，《毛泽东农村调查文集》，第309页。}据不完全统计，1934年上半年，中央苏区的劳动互助社发展到鼎盛时期。其中，瑞金共有社员8987人，兴国县51715人，西江县23774人，长汀县6717人。兴国劳动互助社仅在1934年2～3月间就调剂了7681个人工。\footnote{《和兴国比一比》，《红色中华》第182期，1934年4月30日。}


苏区经济落后，当地农民耕牛奇缺，兴国县永丰区贫农每百户中，“每家一条牛的只有十五家，两家共一牛的四十家，三家共一牛的十家，四家共一牛的五家，无牛的三十家”。\footnote{毛泽东：《兴国调查》，《毛泽东农村调查文集》，第221页。}长冈乡平均百家中有牛25头，“全乡四百三十七家，无牛的约一百零九家”。\footnote{毛泽东：《长冈乡调查》，《毛泽东农村调查文集》，第312页。}正是在此背景下，瑞金县石水乡农民首先创办起犁牛合作社，共同集股买牛，合作管理与使用。合作社的所有耕牛归全体社员所共有，社员大会通过规则进行耕牛管理，选举可靠的人负责饲养耕牛并给予适当报酬，如管理不善，可随时召集社员会议，令其赔偿损失。犁牛合作社以最小的成本使农业生产必不可缺的耕牛资源发挥出最大的效用。


中共在农业合作上的努力体现出其认真做事的劲头和卓越的组织能力，这确实是能够发挥其长处的领域。苏维埃时期，政权对社会的组织和控制达到了中国历史上从未有过的程度，大部分的经济活动被纳入到党和政府的指导之中，这使苏维埃具有很强的控制力，对于解决劳动互助、农具和耕牛合作这样的具体问题，应该是举重若轻，为战争资源的汲取也作出了巨大贡献。相对而言，涉及商业经营领域的合作机构，处理起来就不像在农业合作领域那样得心应手。


营业性的合作社，原则上应为自愿组合的经济实体，毛泽东记载的长冈乡合作社运作情况应为其中的范例：

\begin{quote}

村社、乡社时，社员及红属买货，每千文减五十，即百分之五。非社员不减，但照市价实际上便宜些，一串钱货便宜二十文上下，即百分之二……区社去年九月至今年三月（半年），四百多元本钱赚了六百多元，以百分之五十为公积金，百分之十为营业者及管理委员、审查委员的奖励金，百分之十为文化教育费（为俱乐部、学校及红属儿童买纸笔），百分之三十分红。为了增发红利，鼓励社员，临时将教育费取消（以后应该恢复），共分红百分之四十，每人分了一串钱。分红时，清算账目，悬榜公告。\footnote{毛泽东：《长冈乡调查》，《毛泽东农村调查文集》，第316页。}
\end{quote}


遗憾的是，能够达到长冈乡这样运作状态的并不多。在合作社的发展过程中，像苏维埃这样拥有强大控制力且承担着为社会主义准备发展条件责任的政权，面对着商业精神和经营能力均十分缺乏的社会现实，政府推动之手仍然非常明显：

\begin{quote}

有许多合作社向来不举行社员大会，他们未曾以社员为合作社的主体。许多地方合作社的同志都回答我们，认区委区苏是他们的上级机关，因此，像汀州市粮食合作社就根本忘记了自己的社员，合作社的监察委员会竟不由社员所选举，而由风马牛不相关的其它组织（当然在某种意义上），如工会、政府、列宁书局等委派代表来组织。\footnote{寿昌：《关于合作社》，《斗争》第18期，1933年7月15日。}
\end{quote}


政治力的推动的确促进了合作社的发展，表现出高效率的一面，但滋生的问题也不容忽视。按照当时中共领导人的解释，合作社“应该是在经济战线上反对投机商人与富农的剥削，打破敌人封锁的生力军”，\footnote{亮平：《目前苏维埃合作运动的状况和我们的任务》，《斗争》第56期，1934年4月21日。}是阶级斗争在经济战线上的反映。因此，中共对合作社给予了资源和政策上的大力支持，尤其是消费合作社，被赋予“便利工农群众贱价购买日常所用之必需品，以抵制投机商人之操纵”\footnote{《中华苏维埃临时中央政府关于合作社暂行组织条例的决议（1932年4月12日）》，《革命根据地经济史料选编》（上），第87页。}的任务，其结果，在合作社发展较好地区，实际上取代商人成为农村商品供给的主要来源。以兴国为例，1933～1934年，“全县合作社的营业（包括进出货），十一十二两月为十一万二千元，一二两月为十二万二千元，这就是说，每月平均有三万元以上的商品，是经过了消费合作社的系统，供给兴国群众的消费者”。\footnote{亮平：《目前苏维埃合作运动的状况和我们的任务》，《斗争》第56期，1934年4月21日。}对于当时只剩17万人口的兴国而言，每月3万元的商品供应，全年即为36万元，合作社人均供应两元多商品。根据1930年代的调查，江西新淦谦益村870余人，每年的农业用品和生活用品购买总量约3200余元，\footnote{《视察丰城清江新淦三县报告书》，《军政旬刊》第5期，1933年11月30日。}每人均摊约3.68元。考虑到谦益村地处平原，未受到战争严重影响，生活相对安定，基数要高一些，兴国两元多的商品供应意味着合作社几乎已经包揽了该县的商品供应。\footnote{当时，中央苏区农村主要的生活消费是盐、布匹及少量日用品，像猪肉、鸡鸭、食用油、茶叶等一般都是自给或小范围的圩场流通，人均消费大概也就在一元左右。相对而言，生产用品需求稍高一些。}事实上，毛泽东在上杭才溪乡的调查就发现，这里“卖‘外货’的私人商店，除一家江西人开的药店外，全区绝迹（逐渐削弱至此），只圩日有个把子私人卖盐的，但土产如豆腐等，私人卖的还有”。\footnote{毛泽东：《才溪乡调查》，《毛泽东农村调查文集》，第346页。}所以，正如当时报告所指出的，中央苏区的合作社“极大多数不是发展苏区经济，便利工农群众，而是一部分群众集股的商店，大多数是政府没收的店子或出资办的。实际上都是垄断市场图谋赚钱，根本违犯了合作社组织的原则与作用”。\footnote{《江西省第一次工农兵苏维埃大会——财政与经济问题的决议案》，《中央革命根据地工商税收史料选编》，第86页。}


合作社在商品流通中的控制地位，固然达到了打击所谓投机商人赚取利润，保护资源不使外流的目的，但当合作社取得垄断地位后，一个没有竞争的市场使得合作社掌握着绝对的定价权，其追逐利润的投机程度甚至超过原来的商人。当时的文章披露：“全总的合作社，以二角钱的股本在半年余得三倍的利润，瑞金县总社以一万元资本在三月内获利五千元，兴国合作社每月能获利五千元以上。”\footnote{亮平：《目前苏维埃合作运动的状况和我们的任务》，《斗争》第56期，1934年4月21日。}兴国上社区合作社从1932年9月至1933年3月半年时间内，四百多元本钱赚了六百多元。\footnote{毛泽东：《长冈乡调查》，《毛泽东农村调查文集》，第316页。}这样惊人的盈余是一般的商业利润不可想象的。当时《斗争》发表的文章批评：“合作社的商品价格，常时随商人价格而高低，而且有些合作社，竟把价格规定得与市上相差无几，甚至相等。”\footnote{亮平：《目前苏维埃合作运动的状况和我们的任务》，《斗争》第56期，1934年4月21日。}如果仅仅只是这样的话，无法解释合作社如何能够获取这样的暴利，在一般商业利润难以超过10\%情况下，即使其商品流通十分畅捷，也难以想象半年可以获得三倍的利润。有理由认为，合作社在某些货品上利用其垄断地位获取了超额利益。


垄断获取的超额利润，在监督机制不健全的背景下，极易带来贪污浪费，这就是当时文件所反映的：“合作社的组织，多半是由党与苏维埃、群众团体的负责人所创办、所主持的；合作社的营业，不是为了适应社员的要求，而是为了赚钱。特别是一些工作人员，借着合作社机关的招牌，大做投机生意，垄断市场，贩卖谷盐进出口，就成为整个合作社的主要营业。”\footnote{项英：《于都检举的情形和经过》，《红色中华》第168期，1934年3月29日。}政治力量的过度介入，使有些地区合作社变成了政府机关的私产，为干部上下其手提供了空间。当时材料披露了一批贪污浪费实例：博生县合作总社，“有严重的贪污。主要的是营业部长罗科培，在旧年11月至2月12日止查日记簿，总数上少收九十六元八角七分，帐目计角改元贪污一百四十六元八角三分。在杂用上付出二十八元，作伙食报销，总计贪污三百余元”。\footnote{《轻骑队的活动在博生》，《青年实话》第3卷第19号，1934年4月15日。}西江砂星区高屋乡合作社“赊欠、挪扯得一塌糊涂”。\footnote{《挪扯合作社俱乐部公款》，《红色中华》第142期，1934年1月10日。}赣县青溪、山溪两区合作社更加荒唐，他们“觉得做生意赚钱难得费，于是装运谷子出口，不购买半点群众需要的日常货物，就完全买现洋拿回来做钱生意，他们将现洋做四吊八百文或五吊钱来收买苏维埃国币”，\footnote{《破坏国币的合作社》，《红色中华》第131期，1933年12月2日。}干起了倒卖货币的勾当。亏损现象在合作社中也十分普遍，由于贪污和经营的问题，“踏迳区落村乡的消费合作社，以前有五千毛小洋资本的，现在亏得只剩一千多毛，沙心区的消费合作社，以前有四千毛资本，现在也蚀去一半以上，樟脑合作社的一百余元资本完全蚀光”。\footnote{亮平：《经济建设的初步总结》，《斗争》第29期，1933年10月7日。}正因不少合作社难以取信于民，当时苏维埃政府在发动群众加入合作社时会这样告诉民众：“加入合作社的好处是很多的，你们不要怕经手人作弊，我们可以在社员中选举三个最公直而且非常精细的人来组织审查委员会，专门监督与检查管理委员会的工作。如果有贪污弄弊事情，可以报告苏维埃政府追缴并惩办他。”\footnote{《想激进改善自己的生活，只有迅速加入消费合作社》，《湘鄂赣革命根据地文献资料》第2辑，第247页。}侧面反映出民众对合作社公正性的疑虑。由于合作社存在的如上问题，其兴衰趋势大致如下所述：

\begin{quote}

合作社开始的时候，并没有广泛的宣传工作，多半是由几个负责人发起，但群众的加入难保没有一部分勉强（或因沿户收费，或因买不到盐，作些希望）；但每一个合作社开办以后对群众的需要品确有相当的调解作用，比一般商店价较为低廉（但有时又较贵），但不到一季，就本干利空而塌台。主要原因，就是合作社开设于某处是成了那个地方党与苏维埃负责人的赊拿所，每次货来了，一般的现象是党和苏维埃负责人先得到东西，有些加入了合作社的群众买不到东西，甚至丢却工夫连跑几天仍买不到东西……这就是合作社一般的现象，也就是一般的结果。\footnote{《关于湘鄂西具体情形的报告（1932年12月》，《湘鄂西苏区革命历史文件汇集》第1册，第285页。}
\end{quote}


中央苏区在经济力的组织上，有成功的经验，也存在许多有待进一步探索和解决的问题，面对当时资源有限的客观现实，中共的努力极大地发挥了社会经济的潜力，初步显现出中共处理社会经济问题的理念和能力，但由于中央苏区本身资源的限制、国民党方面的封锁压迫，长期战争的巨大消耗及中共在经济建设过程中显现的某些弱点，使其终究无法支持持久战争的消耗。作为中共的首次建政尝试，在经验有待丰富尤其是战争状态的背景下，遭遇挑战实属正常，其中的某些问题，事实上在中共日后长期的革命乃至执政道路上，仍然是有待深入认识和克服的难题。

\section{财政紧张下的民众负担}


资源匮乏给经济建设带来巨大难题同时，更造成财政严重紧张。苏维埃初期，在苏区内乃至赤白边境地区打土豪，既是红军的政治任务，也为中共解决财政问题提供了便捷渠道，红军常有“筹款部队专门负筹款责任”。\footnote{《江西省苏扩大的第二次全体执委会——关于财政问题决议案》，《中央革命根据地工商税收史料选编》，第186页。}经由此，中央苏区财政状况一度相当不错，甚至可以给上海的中共中央提供财政支持。1931年2月，共产国际远东局的报告谈到中共情况时说：“工作将会进行得更快些，因为现在有钱。中央从毛泽东那里得到了价值约10万墨西哥元的黄金，从贺龙那里得到了1.5万墨西哥元的黄金。”\footnote{《共产国际执行委员会远东局给共产国际的信》，《共产国际、联共（布）与中国革命档案资料丛书》第10册，第143页。}1932年1月，赣东北苏区在自身财政已遇到困难时，仍响应中共中央“各苏区要将筹款帮助中央作为目前战斗任务之一”的指示，先后给中共中央送去纯金条350两。\footnote{《中共赣东北省委关于财政问题向中央的报告（1932年1月）》，《闽浙皖赣革命根据地》（上），第408页。}不过，这种非常态的收入毕竟不具可持续性，虽然打土豪一直是苏区收入中的重要部分，1934年江西省苏维埃还报告：“今年一月份以来，四十天内，全中央区筹款二十五万元。”\footnote{《江西省苏扩大的第二次全体执委会——关于财政问题决议案》，《中央革命根据地工商税收史料选编》，第186页。}但由于国民党军的封锁和压迫，打土豪越来越困难，其在苏区财政中比重呈明显的下降趋势。


在红军壮大时期，攻打财富集中的城市是苏区缓解财政困难的重要手段。1932年5月，周恩来等在解释攻打漳州的行动时谈道：

\begin{quote}

由于缺乏资金，我们又决定扩大北部的苏区。但资金还是不够。我们又改变了先前的决定，决定派一个军团去福建，以解决资金问题。

在朋外的第3军团始终未能解决资金问题，因此……开始向漳州进攻。在漳州募集资金后，我军准备回过头来进攻广东来犯福建和江西之敌。\footnote{《周恩来、王稼祥、任弼时和朱德给中共中央的电报（1932年5月3日）》，《共产国际、联共（布）与中国革命档案资料丛书》第13册，第147页。}
\end{quote}


李德在1933年初的报告中则提到红军攻打城市，“直到商品和钱款被运出”。他解释道：“中央军团甚至得到中央苏维埃政府的专门指示，进行这种实际上的游击行动。面对我们苏区经济基础的局限和敌人的严密封锁，这常常是获取以后进行战争所需物资的唯一出路。”\footnote{《布劳恩关于中央苏区军事形势的书面报告（1933年3月5日）》，《共产国际、联共（布）与中国革命档案资料丛书》第13册，第337页。}然而，随着国民党军战线紧密，实力厚增，这一方法到苏区时代后期事实上不再可行。


通过打击对立面以得到财政支持做法难以为继后，为保持政权的顺利运转，财政对民众的依赖不得不日渐增强。1932年，可以获得较为确切资料的几个苏区县的征税状况是：永丰9个区102711人规定的土地税为46396元，\footnote{《中共永丰县委两个月（十月二十日至十二月二十日）冲锋工作报告》，《江西革命历史文件汇集（1932年）》（二），第455～457页。该书中对人口及税额总数计算有误，笔者根据其提供的基础数字进行了重新计算，同时由于八都区未提供税额，计算中剔除了八都人口。}人均0.452元。于都全县土地税额为人均77000多元，\footnote{《中共于都中心县委两个月冲锋工作报告（十月二十日到十二月二十日）》，《江西革命历史文件汇集（1932年）》（二），第482页。}该县当时人口约数19.1万，人均0.403元。胜利县人口15.33余万，胜利县土地税到当年12月已收款42416元，未收款11233元，另收税谷1300多担，\footnote{《中共胜利县委两个月冲锋工作报告（十月二十日到十二月二十日）》，《江西革命历史文件汇集（1932年）》（二），第473～474页。原统计计算有误，笔者根据基础数字进行了重新计算。}以每担谷5元计，谷价6500元，总计该县税额为60149元，人均0.392元。会昌1932年12月收到当年土地税57691元，占十分之六强。\footnote{《中共会昌县委两个月（十月二十日到十二月二十日）冲锋工作总报告》，《江西革命历史文件汇集（1932年）》（二），第428页。}以此计算当年该县土地税额应为95000元左右，除以该县215000余的人口数，人均0.442元。从这几个有相对确切数据的县份看，1932年苏区土地税额大致为人均0.4～0.45元。和苏区之外国民政府控制区域比，单纯从田赋看，苏区征收比例并不低，不过，苏区在土地税之外，其他负担较轻，1933年，建宁每月收取营业税、烟酒屠宰税、进出口税共计800元，租款480元，以年计共15360元，\footnote{《闽赣省财政部七、八、九三个月工作和八、九、十、十一、十二、一、二七个月筹款计划》，《闽赣苏区文件资料选编》，第58～59页。}按建宁5万人左右人口推算，人均合0.3元。两者相加，就是苏区人民的所有常规负担，由于无须承担各种各样的地方附加和摊派，从总量上看，和国民政府区域人均1元以上的税赋负担比，苏区人民负担有较大减轻。


然而，随着军事紧张、资源消耗加剧，苏区人民在资源紧缺背景下，为支持战争、维护苏区不得不承受更大压力，付出巨大牺牲。这一局面的出现不应简单视作政策错误，而是客观环境使然。1932年6月、10月和1933年7月，由于财政紧张，苏维埃中央政府先后三次发行公债60万、120万、300万元。第一期公债大都用抵交土地税方式陆续归还，第二、第三期则基本成为无偿的贡献。如果把1933年发行的300万元公债平摊到苏区约300万人口中，人均负担将有成倍的增加。1934年苏区政权借谷近百万担，每担以价值5元算，总值近500万元，而此时中央苏区人口已下降到200万人左右，人均实际担负两元多，加上其他支出，苏区民众人均负担已超过3元。


另外，苏区民众常常自发或有组织地慰劳军队。湘赣苏区1933年慰劳军队的主要物品包括：鞋15550余双，棉衣798件，草鞋28911双，米1294石，各种菜品15万多斤，猪肉4544斤。\footnote{赵可师：《赣西收复区各县考察记》（四），《江西教育旬刊》第10卷第8期，1934年8月11日。}宁都1932年7～9月慰劳红军物品包括各种鞋11081双，猪肉183斤，鸡鸭35只，蛋82斤，鱼28斤，果饼18担，花生15担，蔬菜40担。其他还有水果、牙粉、牙刷、纸烟、银洋等。\footnote{《中共宁都县委七八九三个月工作报告》，《江西革命历史文件汇集（1932年）》（二），第127～129页。}胜利县1932年8月的慰劳品是：布草鞋3200双，布套鞋920双，头牲143斤，鸡蛋5730个，生猪5只，牙粉103包，糕饼557付，梨子四担等。\footnote{《中共胜利县委工作报告》，《江西革命历史文件汇集（1932年）》（二），第190～191页。}这些慰劳品由于不是定期定量供应，所以难以得出一个具体的供应数据，但长年累月算起来，也是一笔不小的开支。而且，随着战争日益紧张和兵员的不断扩大，这样的支出还呈日渐增多的趋势。


客观而言，对于局限固有地区的中央苏区而言，要维持政权的运转和保卫苏区，财政压力不断加重是一个难以避免的结局，解决这一问题，很难有其他的疏浚渠道，只有依赖群众的帮助，而这必然会加重群众的负担，并反映到群众的情绪中。虽然大部分群众表现出全力支持中共的热情，但现实的困难仍然使很多人有心无力。1933年中央苏区的财税收入状况就不理想，江西省苏1934年3月的决议中谈道：“去年税收及各种财政收入不能达到如中央财政计划所规定，是革命中重大的损失。”\footnote{《江西省苏扩大的第二次全体执委会——关于财政问题决议案》，《中央革命根据地工商税收史料选编》，第185页。}安远县龙布区苏的会议记录显示，这里的土地税征收遇到困难。1933年3月，该区就要求各乡应把上一年的各种税收在当月23日前收清，但这一要求显然没有实现。4月20日的区苏会议又规定“土地税按期本月底完全收清”。从6月10日该区会议记录看，这一规定仍然没有完成，因为会上再次要求“去年土地税和二期公债票限五天内一律收清算明”。\footnote{《安远县龙布区苏主席团会议记录》，《江西革命历史文件汇集（1933～1934年）》，第53～68页。}苏区中心地区的瑞金沿江区群众也对缴税抵触，抱怨“旧年缴都不要这样重”，“四乡的党团和代表都带群众到区委区苏要求减少”。\footnote{《中共瑞金县委九月份工作综合报告》，《江西革命历史文件汇集（1932年）》（二），第167页。}


战争紧张是造成苏区财政困难的最直接原因，同时，就赣南、闽西这样一个狭小且资源有限地区而言，承担起拥有十多万军队及庞大机关人员的国家机器，实在也是勉为其难。苏维埃组织严密，对基层的控制力达到中国历史上前所未有的高度，与此同时，各级机关十分庞大，对资源的消耗巨大。1931年11月苏维埃临时中央执行委员会第一次全体会议通过《地方苏维埃暂行组织条例》规定，每乡苏维埃政府脱产人员3人，城市（县苏所在地）苏维埃19人，区苏维埃19人，县苏维埃25人，省苏维埃90人，但在实际执行中，一般都会突破此一标准。江西省苏1934年3月公布的各级财政部门人员组成是：省财政部24人，县财政部13～16人，区财政部6～8人，仅县一级财政部门人数就达到规定的县苏总人数一半以上。而且，江西省财政部门还设立了11个关税征收点，每个点都有工作人员9人，检查队一班8～24人。\footnote{《江西省苏扩大的第二次全体执委会关于财政问题决议案》，《中央革命根据地工商税收史料选编》，第188～189页。}总计起来，财政部门的编制数远远高于临时中央政府规定的统一标准，而其实际人数又要突破编制数，当时在财政部门工作的干部回忆：“1933年，我调县财政部当副部长，县财政部有三十多个人，有正副部长、文书，还有会计科（会计七、八个人），税务科（七、八个人），国产管理科（二人），还有伙夫等。”\footnote{《宁都县苏区征收屠宰税、房租税情况——访问彭世鹤记录》，《中央革命根据地工商税收史料选编》，第292页。}县财政部门实际人数已高于整个县苏的规定编制。


财政部门只是整个庞大机构的一个缩影。1933年10月，瑞金县苏维埃工作人员达到302人，远远超过编制规定的25人。区苏维埃一般达到40～60人，以瑞金18区计算，取其平均值即可达900人。加上乡苏维埃人员，瑞金政府工作人员总数应在1500人以上。\footnote{《介绍瑞京裁减间员节省经费的经验》，《红色中华》第169期，1934年3月31日。}瑞金人口20余万，约相当于苏区总人口的1/14，以此推算，苏区县级及其以下政府机关人员可达6万人。加上其他各级机关，干部数量应是一个相当庞大的数字，1934年3月《红色中华》曾刊文号召展开节省运动，其中说道：“政府工作人员每人每日照规定食米量节省二两，以八万人计，每月可节省谷子四千五百石。”\footnote{《为四个月节省八十万元而斗争》，《红色中华》第161期，1934年3月13日。}这里提到的8万人，应是苏区中央对干部数量的概略估计，有相当的可信性。


苏维埃群众组织众多，一些群众组织也有固定编制。如江西省工联会和雇工会固定编制有15人，每月经费165元。\footnote{《江西省职工联合会通告（第二号）》，《江西革命历史文件汇集（1932年）》（一），第2页。}湘鄂赣省苏、省委各机关总人数甚至高达3300余人。\footnote{陈小鹏：《赣西北匪情实况及进剿意见书》，《湘鄂赣革命根据地财政经济史料摘编》（下），湖南人民出版社，1989，第855页。}所以，当时这样的情况并非个别：安远县天心区第四乡“全乡的少先队、儿童团以及什么妇女队等等共四十多人在政府吃饭，弄得该政府忙个不开的办饭主义”。\footnote{翰文：《天心区与各乡苏维埃主席联席会议的经过》，《武库》第10期，1932年2月6日。}同时，各级机关还有一些具体办事人员，如乡级可设伙夫、交通，区、县更多，这些人的津贴并不低于负责人员。虽然总体看，各级机关人员只给不多的津贴，但由于基数庞大，仍是一笔不小的数目。除机关人员外，红军、游击队和一般工作人员也占有相当大的数量，永丰只有十万苏区人口，其脱离生产的地方部队（包括独立团、游击队、模范团、县警卫连）就达933人。\footnote{《中共永丰县委两个月（十月二十日至十二月二十日）冲锋工作报告》，《江西革命历史文件汇集（1932年）》（二），第452页。}


庞大的工作人员队伍，使脱产人员比例空前提高。根据毛泽东对长冈乡、才溪乡的调查，长冈乡全乡1785人，调县以上工作人员34人，如果加上区、乡工作人员，干部人数超过人口比例的2\%。出外工作人员一共有94人，参加红军和游击队者226人，总计出外320人，占人口总数的17.9\%。才溪乡全乡4928人，调外工作的186人，加上参加红军者共1040人，占人口总数的21.2\%。\footnote{毛泽东：《长冈乡调查》、《才溪乡调查》，《毛泽东农村调查文集》，第288、342页。}这样的比例不仅出现在这两个先进区，其他区也大体和其相当。上杭通贤区通贤乡实有劳动力为41人，出外工作人员达到62人，参加红军、做长期夫子和赤少队、模范营者103人，后两者远远多于留在农村的劳动力。\footnote{《各乡劳动出入表（1934年5月27日）》，福建省档案馆藏档435001-91-2-302。}第五次反“围剿”时，中央苏区脱产、半脱产人员总数达到三四十万人，平均每8个人就要负担一个脱产、半脱产人员，民众的负担对象大大增加。张闻天曾经谈道：“常常有这样的同志说，为了中国革命的胜利，农民是不能不牺牲一点的。”\footnote{张闻天：《苏维埃政权下的阶级斗争》，《斗争》第15期，1933年6月15日。}虽然他批评了上述说法，但这种说法相当程度反映着当时的现实。


沉重的负担加剧了苏区财政入不敷出的局面。赣东北苏区1930年11月至1931年3月收入总数不过75000元，但每月支出达7万～8万元；1931年总亏空大洋66511元、金子1205两，\footnote{《中共赣东北省委关于财政问题向中央的报告（1932年1月）》，《闽浙皖赣革命根据地》（上），第406～407页。}严重入不敷出。紧邻中央苏区的湘赣苏区1932年9月至1933年8月底一年的收支账是：大宗收入包括土地税58740元，造币厂盈利50207元，金矿局盈利2426元，营业税20707元，纸业合作社798元，罚没款2855元，豪款10768元，杂收786元，富农捐款5977元，各县缴款7688元，加上其他收入总计161939元。支出包括行政费22227元，帮助各县经费2856元，军费191783元，司法费268元，保卫费13112元，教育费275元，津贴费29249元，国营支出845元，临时费357元，加上其他支出共252612元。\footnote{赵可师：《赣西收复区各县考察记》（四），《江西教育旬刊》第10卷第8期，1934年8月11日。}收支相抵，赤字达9万余元。中央苏区由于机构庞大，军力较强，费用消耗更大。曾任中革军委总动员武装部部长的杨岳彬在投降国民党后谈到1933年度中央苏区的经费使用情况：“截至去年十二月止，每月军事用费（包括匪军伙食及购买药材等）达四十万元以上，各级伪政府经费约十余万元。”\footnote{《匪情实录》，赣粤闽湘鄂剿匪军东路总司令部：《东路月刊》第3、4期合刊。}杨还没有提到各级党组织的费用。以此推算，1933年苏区维持运转的费用达600万元以上。与此同时，收入状况却不容乐观，早在1930年闽西就出现亏空，该年4～10月收入142000余元，支出182000余元，入不敷出4万元。\footnote{《中共闽西特委报告第一号，1930年11月》，《福建革命历史文件汇集》甲8册，第214页。}1932年中央苏区征收的土地税在73.4万元以上。1933年土地税为粮食征收，计22.5万担，按照政府指定粮价5元计算，换算成现金是112.5万元。即使按照市场价格翻番计算，也只有200多万元。维持苏区运转主要必须依靠打土豪及其他非常规收入，同时发行公债弥补亏空。


在中央苏区遭遇财政困难时，苏俄方面曾力图提供支持。中国革命在共产国际和苏俄心目中具有十分重要的地位，看一看1928年共产国际的有关电文，就可以对这一点有清楚的了解。下面是1928年苏俄援助东亚各国的资金清单：

\begin{quote}

拨给中国共产党第二季度（4月、5月和6月）每月12820美元。

拨给日本人共产党每月1025美元。

可以拨给朝鲜人每月256美元作为日常开支。

拨给中国共青团1928年上半年7692美元；可以在这个数目范围内拨给他们。

拨给日本共青团1928年上半年512美元。可以在这个数目范围内拨给他们。

拨给朝鲜共青团1928年上半年460美元。可以在这个数目范围内拨给他们。\footnote{《皮亚特尼茨基给阿尔布列赫特的信（1928年4月3日）》，《共产国际、联共（布）与中国革命档案资料丛书》第7册，第395页。}
\end{quote}


苏俄给中共党人的援助要远远高于日本和朝鲜，这当然是由于中国革命运动开展的程度所决定的，中国是苏俄在东方推动革命的主要期待。在中国革命遭遇困难时，为支持中央苏区的反“围剿”战争，苏俄更是殚精竭虑。1933年10月，共产国际致电中共上海中央局，询问：“请弄清楚并尽快告知，能否购买几架飞机，特别是歼击机。有否希望委托可靠的飞行员把这些飞机从空中提供给苏区？”同时要求上海中央局尽力为中央苏区购买药品和防毒面具，强调：“为达到这些目的，我们可以拨出专项经费。”\footnote{《共产国际执行委员会东方书记处给李竹声的电报（1933年10月12日）》，《共产国际、联共（布）与中国革命档案资料丛书》第13册，第545页。}10月底，由于中共方面与十九路军接触，十九路军答应居间为中共购买武器弹药，共产国际驻华代表致电国际方面，告其“紧急寄出5万元，有购买弹药的可能性”；随后又告以：“购买药品，先急需3万元。”\footnote{《格伯特给皮亚特尼茨基的电报（1933年10月30日）》、《埃韦特给共产国际执行委员会东方书记处的电报（1933年11月10日）》，《共产国际、联共（布）与中国革命档案资料丛书》第13册，第581、617页。}


苏俄和共产国际资助中共的经费，从共产国际代表报告中可见一斑。1933年8、9月份，国际驻华代表埃韦特经手转交中共的经费包括：24.56万法郎、6.16万美元、101452墨西哥元、5000瑞士法郎、1864荷兰盾。\footnote{《格伯特给共产国际执行委员会国际联络部的电报（1934年1月25日）》，《共产国际、联共（布）与中国革命档案资料丛书》第14册，第27页。当时货币含金量为：1934年1月美元贬值前1.50466克，贬值后0.888671克；法郎0.05895克；瑞士法郎0.2032258克；荷兰盾0.334987克（1946年）。墨西哥元当为20世纪初在中国流行的墨西哥鹰洋，比中国银元价值略高。}这应该是一个相当庞大的数目。1933年11月，埃韦特报告其收到的一笔款项甚至达到“300万墨西哥元”。\footnote{《埃韦特给皮亚特尼茨基的电报（1933年11月14日）》，《共产国际、联共（布）与中国革命档案资料丛书》第13册，第618页。}由于此前共产国际驻华代表和中共上海中央局常常向国际叫穷，并声称：“你们寄来的款项，很大比例被我们花掉了。”国际方面也得到消息：“由于无法转给苏区，斯拉文（李竹声——引注）那里存了很多钱。在上海，钱转来转去保存。”\footnote{《埃韦特给皮亚特尼茨基的信》、《皮亚特尼茨基给埃韦特的电报（1934年4月26日）》，《共产国际、联共（布）与中国革命档案资料丛书》第13册，第470页；第14册，第118页。}共产国际希望好钢用在刀刃上，督促上海方面尽力利用钱款为中央苏区提供帮助。1934年5月，共产国际致电中共中央上海局负责人李竹声：

\begin{quote}

通过您收到了来自江西中央关于购买药品、食盐和用于生产子弹的原料的电报。您给他们转寄了我们寄去的用于采购的所有款项吗？如果没有，请马上告知还有多少钱。每月您能无特别风险地给中央转寄出多少。为了利用采购和向江西提供物资的机会，需要在南方，可能的话在澳门设点，并从那里经福建港口建立特殊的联系路线……应该成立一个公司，从事贩卖四川鸦片生意和从四川向江西盗卖白银。这样我们就可以为购买江西红军所必需的东西提供极重要的资金援助。\footnote{《共产国际执行委员会政治书记处政治委员会给李竹声的电报（1934年5月26日）》，《共产国际、联共（布）与中国革命档案资料丛书》第14册，第123页。}
\end{quote}


6月，共产国际直接致电中共中央，指示：“请从苏区和从上海经意大利公司和其它外国公司或者军阀代表处寻找联络途径，以便通过最经济和最可靠的途径购买和提供弹药。你们能否为此建立自己的隐蔽的中介公司？请尝试通过这些公司出售四川红军有的商品，为中央苏区换取武器。”\footnote{《共产国际执行委员会政治书记处政治委员会给中共中央的电报（1934年6月17日）》，《共产国际、联共（布）与中国革命档案资料丛书》第14册，第146页。}


中央苏区撤离计划基本确定后，财政需求更加迫切，7月底，上海方面向共产国际报告：“我们又给苏区寄去5万墨西哥元。到9月中旬还需要寄40万墨西哥元，重复一遍，40万墨西哥元，因为晚些时候，看来几乎没有机会了。”\footnote{《中共上海中央局、盛忠亮和格伯特给皮亚特尼茨基的电报（1934年7月25日）》，《共产国际、联共（布）与中国革命档案资料丛书》第14册，第171页。}从电报透露的数据看，共产国际的支持确实不是一个小的数量，可要将钱和物资寄达中央苏区并不容易，尤其是后者更难完成。9月初，共产国际提出“在中国南方的一个港口建立一个为苏区采购和运输武器、弹药和药品的不大而有效的机构”，\footnote{《共产国际执行委员会政治书记处政治委员会给中共中央的电报（1934年9月4日）》，《共产国际、联共（布）与中国革命档案资料丛书》第14册，第234页。}但由于红军很快撤离，计划根本未及实施。10月14日，当红军已经开始撤离时，王明还在询问“是否还需要在南方建立采购武器的机构？”\footnote{《王明和康生给中共中央的电报（1934年10月14日）》，《共产国际、联共（布）与中国革命档案资料丛书》第14册，第278页。}看来，远水终究难救近火，共产国际和苏俄的帮助，在中央苏区，起到的作用终究有限。

\section{查田运动：理念、策略与现实}


中共领导的苏维埃革命，以土地革命为重要旗帜，因此，苏维埃革命时期，中共在土地问题上投入了大量的精力，具体政策也历经变更，其基本目标均在使普通农民尽可能多地获得土地，以实践中共抑制剥削的阶级革命理念，巩固中共在农村中的群众基础。1933年第五次反“围剿”展开前夕，秉持着这一思路，中共再次在苏区农村展开大规模的查田运动。就文本的宣示看，查田运动旨在于革命战争紧张的形势下，在苏区内彻底清查地主、富农隐瞒成分，进一步在苏区执行“地主不分田，富农分坏田”政策，从而深化苏区内部的阶级观念和阶级斗争，纯洁阶级队伍。同时，由于面临国民党军第五次“围剿”，查田运动当然包含着为即将到来的反“围剿”战争凝聚力量的目标，这也应该是中共中央发动这一运动的初衷之一。


1933年6月，苏维埃中央政府发出《关于查田运动的训令》，接着召开中央苏区八县区以上苏维埃负责人查田运动大会，查田运动在中央苏区迅速集中开展。作为一场阶级革命中的阶级运动，对阶级关系作出判断为其题中应有之义。查田运动以清理阶级关系为发动理由，其对苏区农村阶级关系的判断逻辑上必然是紧张和严重的。运动中下发的文件作出结论，指出苏区虽然经过土地革命，仍然存在着地主富农的强大势力，这些势力的具体体现主要是“那些冒称中农贫农分得土地的地主富农分子”，\footnote{《中央政府通告召集八县区以上苏维埃负责人会议及八县贫农团代表大会》，《红色中华》第85期，1933年6月14日。}需要在运动中加以摧毁。毛泽东也强调：“查田运动是一个剧烈的残酷的阶级斗争，必须发动最广大群众热烈起来参加斗争形成群众运动，才能保障阶级路线的正确执行，才能达到消灭封建残余势力的目的。一切脱离群众的官僚主义命令主义工作方式，是查田运动最大的敌人。查田运动的群众工作，主要是讲阶级，通过阶级，没收分配，及对工会贫农团的正确领导等。”\footnote{毛泽东：《查田运动的群众工作》，《斗争》第32期，1933年10月28日。}


作为运动名义上的领导者，毛泽东这时处境微妙。由于中共中央机关的到来，毛泽东实际已不参加重大事务的决策，他此前的工作也被新的领导层摆在放大镜下加以检验。1933年初，中共中央机关刚刚到达苏区，就强调要加紧推进查田，为此，1933年2月，苏维埃中央政府要求：“田未分好，或分得不好的地方……要马上发动群众，重新分田。”\footnote{《中华苏维埃共和国临时中央政府土地人民委员部训令》，《红色中华》第52期，1933年2月13日。}同时，土地部组成工作组，开始在瑞金云集区等地开展查田试点。这样的举动，多多少少体现出一种不信任的态度，毛泽东对此自然心知肚明。因此6月1日毛泽东解释开展查田运动的原因时指出，这是由于苏区内部斗争发展的不平衡，导致一些落后地区“远远落在先进区域之后”，这种地方“占了中央区差不多占百分之八十的面积，群众在二百万以上”。\footnote{毛泽东：《查田运动是广大区域内的中心重大任务》，《论查田运动》，晋察冀中央局翻印，1947，第4页。}强调先进和落后地区的差异，潜台词是要表明各地存在的问题主要在于执行的偏差。但是苏区中央局次日发布的决议则批评道：

\begin{quote}

党和苏维埃政权过去对于土地问题解决的不正确路线（如“抽多补少，抽肥补瘦”，“小地主的土地不没收”等），在许多区域中，土地问题还没有得到彻底的解决。有些区域中虽然已经分配了土地，但是地主豪绅与富农常常利用各种方法（或者假装革命混入党苏维埃机关，或者利用氏族的关系和影响，或者隐瞒田地，或者以物质的收买，政治的欺骗，武力的威吓），来阻止雇农贫农的积极性的发展，以便利他们的土地占有，甚至窃取土地革命的果实。\footnote{《中央局关于查田运动的决议（1932年6月2日）》，《中央革命根据地史料选编》（下），第480页。}
\end{quote}


中央局丝毫不留情面，将过去的土地革命路线定性为不正确的路线，这样的批评不可不谓严厉，在初来乍到的中共领导人看来，此前苏区执行的土地政策是所谓“富农路线”，这也成为他们不点名批评毛泽东的重要理由。然而，在1929年共产国际指责中共六大对富农让步后，各苏区基本都执行“地主不分田，富农分坏田”政策，中央苏区也不例外，在这样的原则问题上，当年中共的政治生态下很难会有别的选择。毛泽东之所以遭受批评，关键不在于他对富农的态度，而是源于其对苏区环境下土地革命极有可能触及中农利益的担忧。在查田运动的动员报告中，他明确提道：

\begin{quote}

联合中农，是土地革命中最中心的策略，中农的向背，关系土地革命的成败。所以要反复向群众说明这个策略，说明侵犯中农利益的绝对不许可的。为了联合中农不侵犯中农利益起见，要提出“富裕中农”来说明它，要着重说明富农与中农交界地方，使富裕中农稳定起来。\footnote{毛泽东：《在八县查田运动大会上的报告》，《论查田运动》，第6页。}
\end{quote}


毛泽东如此重视中农和富农的界限，煞费苦心地在富农和中农中间提出富裕中农的概念，防止混淆中、富农，应该有他深思熟虑的想法。如前所说，中央苏区土地占有比较分散，阶级分化不甚明显，当中共开展阶级革命时，地主、富农的有限资财往往很难满足普通农民改善生活的愿望，在均平的旗帜下，生活高过平均水平的中农很容易成为平均的对象。数年的土地革命实践证明，当打击地主、富农时，中农是最容易被误伤的对象，而中农作为农村最具实力的一个阶层，对它的错误打击常常牵一发而动全身。


1933年6月，苏维埃中央政府发出《关于查田运动的训令》，接着召开中央苏区八县区以上苏维埃负责人查田运动大会，查田运动在中央苏区迅速集中开展。作为一场阶级革命中的阶级运动，对阶级关系作出判断为其题中应有之义。查田运动以清理阶级关系为发动理由，其对苏区农村阶级关系逻辑上必然是紧张和严重的，运动中下发的文件作出结论，指出苏区虽然经过土地革命，但仍然存在着地主富农的强大势力，而这些势力的具体体现又主要是“那些冒称中农贫农分得土地的地主富农分子”，\footnote{《中央政府通告召集八县区以上苏维埃负责人会议及八县贫农团代表大会》，《红色中华》第85期，1933年6月14日。}需要在运动中加以摧毁。毛泽东也强调：“查田运动是一个剧烈的残酷的阶级斗争，必须发动最广大群众热烈起来参加斗争形成群众运动，才能保障阶级路线的正确执行，才能达到消灭封建残余势力的目的。一切脱离群众的官僚主义命令主义工作方式，是查田运动最大的敌人。查田运动的群众工作，主要是讲阶级，通过阶级，没收分配，及对工会贫农团的正确领导等。”\footnote{毛泽东：《查田运动的群众工作》，《斗争》第32期，1933年10月28日。}


事实上，查田运动推开后，打击中农立即成为现实。查田运动前，经过数年革命的中央苏区对地主、富农的清查是比较彻底的。当时苏区拥有较多土地的地主、富农占总人口数的7\%～8\%，查田运动前中央苏区清查出来的地主、富农人口数占到总人口的7\%左右，\footnote{1933年，会昌全县总人口206866人，其中地主、富农人口13828人，占总人口近7\%（《会昌查田运动进行情况》，《红色中华》第106期，1933年8月31日）；胜利县总人口10万人左右，地主富农有1454家，约合7000人左右，也占7\%。参见王观澜《胜利县继续开展查田运动经验》，《斗争》第61期，1934年5月26日。}这一比例和后来得出的全国范围8\%左右比，尚属正常，考虑到中央苏区地主富农经济不发达，其实这一数字本身或许已不无偏高。在此背景下，中共中央领导人主观认定苏区农村中还存在未发现的大量地主、富农，要求深挖隐藏的地主、富农，各地为完成中央要求不得不尽力寻找靶子，以避免被扣上机会主义、动摇妥协的帽子，这就不可避免地使许多贫中农尤其是中农成为所谓“隐藏的地主、富农分子”。而且查田运动大规模铺开后，相应的阶级划分这样一个十分关键的配套政策却没有跟上。在缺乏具体标准时，普通农民区分成分最直观的感受就是生活水平的高低，这种认识在中共各级干部中间也或多或少存在，因此，农村相对生活较好的中农极易成为打击对象。瑞金踏迳区采取普遍清查的办法，“查得一部分中农恐慌逃跑，躲到山上”。\footnote{毛泽东：《查田运动的初步总结》，《斗争》第24期，1933年8月29日。}“有的地方普遍查田，甚至有专门查中农的，说中农中最容易躲藏富农（如会昌的某处），瑞金每个区都发生把中农或富裕中农当富农打的事情……有三个区发生中农上山。”\footnote{《中央关于查田运动的第二次决议》，《中共中央文件选集》第10卷，第337页。}


作为阶级斗争的一个部分，查田运动具有强烈的政治色彩，主观认定以政治正确的姿态强行进入现实，使实际执行者几乎没有讨价还价的余地。在政治正确气氛影响下，很多干部为避免犯错误，主持评定成分时多就高不就低：“查田查阶级方式，往往是开大会，提出名单来问群众某人是否地富、某人是否‘AB团’，叫群众举手。如果群众不举手，便说群众与反革命或地富妥协，于是群众害怕，只好大家举手。因此群众很怕我们，离开我们。”\footnote{李六如：《谈湘赣苏区土地革命》，《回忆湘赣苏区》，江西人民出版社，1986，第89页。}闽西甚至“有因争论阶级而枪毙贫农的事件发生”。\footnote{《中共福建省委工作报告大纲（1933年10月26日）》，福建省档案馆、广东省档案馆：《闽粤赣边区革命历史档案汇编》第1辑，档案出版社，1987，第387页。}而贫农为在运动中获取利益，也愿意将生活较好者定高成分，有时，群众的主张甚至会局部主导运动的气氛。


在多种因素共同作用下，各地高定成分现象相当突出，不仅是中农，甚至贫农、工人也被作为打击对象：“一人在革命前若干年甚至十几年请过长工的，也把他当作富农”；“把稍为放点债，收点租，而大部分靠出卖劳动力为一家生活来源的工人当地主打”；“建宁的城市、里心、安仁等区，共计中农、贫农被误打成土豪的有五十余家，还有一个工人被打成土豪的”。\footnote{毛泽东：《查田运动的初步总结》，《斗争》第24期，1933年8月29日；刘少奇：《农业工会十二县查田大会总结》，《斗争》第34期，1933年11月12日；邵式平：《闽赣省查田突击运动的总结》，《红色中华》第181期，1934年4月28日。}任职于保卫局的童小鹏回忆，其“家庭出身是贫农，但一九三三年因‘左’倾错误政策错打成地主、富农，后虽经纠正，但父兄因流离而死亡”。\footnote{童小鹏：《军中日记》，解放军出版社，1986，第219页。}随意拔高成分的情状，正如于都河丰区委组织部长抱怨的：“工作团是私打地主，将来有一碗吃的人都会被打为地主的。”\footnote{项英：《于都检举的情形和经过》，《红色中华》第168期，1934年3月29日。}


查田运动开展后的三个月内，中央苏区查出所谓的漏划地主6988户、富农6638户。其中，瑞金6～7月间查出地主608户，富农669户，收回土地60591担。石城共查出地主94户、富农179户，收回冒充中农、贫农的地主、富农的土地11200多亩；查出混进苏维埃政权机关的所谓地主23人、富农31人、反革命5人。\footnote{揭国发、刘化尧：《查田中的斗争》，《石城文史资料》第4辑，政协石城文史资料委员会，1991，第15页。}这些被中共中央领导人作为查田运动的成果，也由此证明查田运动的必要性。但是细细检证这些成果，却未必那么经得起推敲。


根据毛泽东30年代初在中央苏区所做多次调查，当时人均拥有6～10担谷田仅为够吃的标准，\footnote{参见毛泽东《兴国调查》、《木口村调查》等。《木口村调查》所列几位中农均处于这一范围内。（《毛泽东农村调查文集》，第283～284页）当时一担田（3～4担田合一亩）约能收谷百斤左右，考虑到家用消耗及留种等因素，800斤谷为基本的温饱标准。}实际生活相当于中农的水平。但在查田运动中，人均拥有9担田、每年需租进田地耕种的家庭也被定为地主。\footnote{《中央土地人民委员部1933年7月13日给瑞金黄柏区苏的一封信》，《红色中华》第95期，1933年7月23日。}以1933年7、8、9三个月中央苏区查出地主6988家、富农6638家，收回土地317539担计，\footnote{毛泽东：《中华苏维埃共和国中央执行委员会与人民委员会对第二次全国苏维埃代表大会的报告》，《中央革命根据地史料选编》（下），第320页。}由于实行“地主不分田，富农分坏田”政策，这实际即为新查出地主的全部土地及富农多余土地之和。以两个数字相衡量，考虑到其他因素，这些所谓的地主每家拥有土地平均约30余担，人均拥地不足10担，仅在温饱线上。同一时期，公略县查出地主381家（家庭总人口1181人），没收地主土地5168担，\footnote{《公略查田运动的检阅》，《红色中华》第125期，1933年11月14日。}新查出的地主家庭每人平均占地也只有10担左右。在这些人中，不乏丧失劳力或因各种原因被迫请人耕种，结果被定为地主者，甚至有红军家属因请人耕种而被定为地主者。福建汀州在查田运动中就有没收“工人、雇农、红军家属财产”\footnote{《福建查田的经验与教训》，《红色中华》第114期，1933年9月30日。}的情况。有些地区确定成分时，“拿剥削的种数，去分别地主与富农的成份。三种剥削的叫做地主，两种剥削的叫做富农。比如请了长工，收租，又放了债，则不管他家里有几人劳动，总之他就是地主了。”\footnote{毛泽东：《查田运动的初步总结》，《斗争》第24期，1933年8月29日。}湘赣省“有一个贫农，查成分查了人家七代，结果被错划为地主”。\footnote{王首道：《回忆湘赣苏区》，《湘赣革命根据地》（下），第848页。}


衡诸常理，以中共强大的组织力，加上武装力量的直接推动和介入，在苏维埃革命展开数年后，苏区仍然会存在相当强大的地主富农势力，多少让人觉得有些不可思议。但对于陆续进入苏区的中共中央领导人而言，这样的判断并不突兀，作为年轻的革命理想主义者，苏区人民的现实状况难以使他们满意。虽然贫苦农民分得了土地，但由于革命、战争、灾害及农村社会客观状况多方面的因素，1930～1932年，苏区生产出现下降局面，农民贫穷状况未得到根本改变。这显然是相信革命可以立竿见影的理想主义者们难以接受的，因此，认为地主豪绅“窃取土地革命的果实”，就其思考逻辑而言确也顺理成章，而且，解决的办法必然是要加紧对地主、富农的打击，“使得土地革命的利益完全落在雇农贫农中农身上”。\footnote{《苏区中央局关于查田运动的决议》，《中共中央文件选集》第9卷，第207页；博古：《为着布尔什维克的春耕而斗争》，《红色中华》第51期，1933年2月10日。}值得注意的是，这种判断并不是在查田运动中最后一次出现，抗战胜利后中共重新领导开展土地革命时，在许多地区这样的判断又一次集中显现。查三代、查“封建尾巴”，打击中农，许多现象几乎就是1930年代历史的重演，韩丁即认为“可以和三十年代的过火行为相比”，而这又恰恰发生在中共集中批判苏维埃革命时期“左”的错误仅仅数年之后。韩丁以自己在潞城张庄的实际经历生动地描述了这种观念的根源：

\begin{quote}

这样的一场运动是以两项假设为前提的。第一、潞城县仍然存在严重的封建剥削；第二、一大部分农民仍然没有翻身。

既然绝大多数有钱有势的人家已经丧失了公开的财产和一部分地财，那么第一项假设就难以成立了。然而，乡村里还有数以千计的贫苦农民很缺乏成为独立生产者所必需的生产、生活资料，因此财产没收似乎还不够彻底。实际上，持久难除的贫困成了那两项假设的根据，成了进行一场新斗争的理由。穷人真要翻身吗？就得再多找出一些财产来。\footnote{〔美〕韩丁：《翻身》，北京出版社，1980，第146、229页。}
\end{quote}


对于事后的观史者而言，注意到查田运动的上述因素，不仅可以发现查田运动本身的复杂性，而且也许可以为认识当年屡屡出现的“左”的错误找到有益的路径。不过当时的中共中央领导人并不可能有这样的反省机会和时间，当问题被抬到阶级对抗的高度时，足以刺激他们绷紧的弦，由于此，查田运动的发起和紧张化几乎是不可避免。即使是熟谙苏区状况并领导过此前土地运动的毛泽东，亲身领导运动之初，虽然以自己对苏区的了解，极力在其中留下余地，但对查田的原则也并没有提出异议。吊诡的是，查田运动是建立在此前革命不彻底的判断基础上，很容易让人联想到是对毛泽东此前主持的土地革命的变相批评，而毛泽东却在名义上担负起领导运动的任务，这其中，既有服从组织的因素，也或有判断上的困难，和军事方面的坚持比，显然，涉及这样敏感复杂、一时间难以理清头绪的政治原则问题，毛泽东的态度要谨慎得多。


更让我们体会到问题复杂性的还在于，当时中共中央在查田运动中屡屡提到的农村阶级关系尚不分明的说法，其基础虽是出于教条化的理论推导，但在苏区的现实环境中，又不是毫无脉络可寻。在赣南、闽西土地集中程度有限，宗族组织具有强大影响力背景下，公式化的农村阶级分化并不一定是当地农民实际状况的真实写照，而以这样的观念指导的运动遭遇反激也势所必然。黎川梅源“因系吴氏一姓，血统亲密之故，极能互为维护，毫无幸灾乐祸、趁火打劫心理，复无有广大土地之大地主，以为自相仇杀之导火线”，因此，土地革命展开后，虽然该地也进行土地平分，但具体运作中多有问题。分田时，农民常会用自己的方式阳奉阴违：“将能见到者，以插标为记，各分谷田五石，其余陇亩不能见到之者，秘不均分。其已分配者，耕作之后，仍将其所收之谷，按佃户例，送还原主。土匪因此怀疑分配不实，于是一再举行分配，并有所谓查田运动之新花样出现，但举行结果，依旧如故。”\footnote{《黎川梅源概况》，《汗血月刊》第1卷第4期，1934年4月20日。}在赣南、闽西许多地区实际社会状况与此大同小异背景下，这样的现象当非个案。正因如此，在国民党重新占领赣南后，杨永泰注意到：

\begin{quote}

广昌和黎川，都是分过田的。分田的小册子和分田的标帜，统统都有，但是田畦还是维持着原状，并未敢加以破坏。因为人民反对破坏，说是：田畦一经破坏，田里就不能蓄水，没有水就没有办法耕种，共党也无法可施，只好迁就作罢。\footnote{杨永泰：《革命先革心，变政先变俗》，《新生活周刊》第1卷第15期，1934年8月6日。}
\end{quote}


杨的说法从中共自身文件中也可得到证实，中共中央的有关文件谈到相对多的土地拥有者“利用氏族的关系和影响”\footnote{《苏区中央局关于查田运动的决议》，《中共中央文件选集》第9册，第207页。}保存土地的现象；而《红色中华》等舆论机关也披露过诸如石城县苏维埃主席邓海如以地方家族观念庇护同姓同村的地主、富农的事例。\footnote{《石城八月份查田总结》，《红色中华》第111期，1933年9月21日。}这些干部包庇事例的出现，除宗族、地方观念因素外，相当部分地区土地拥有者和其他农民无论在经济、政治还是社会生活上并没有真正的鸿沟也是一个不容忽视的因素，因此，即使是进入主流阶层的干部，在尚未完全自觉坚持中共阶级分析立场时，其长期耳濡目染的社会现实仍然影响着他们的选择。事实上，在中央苏区的部分地区，土地革命的开展相对顺利，而在另一部分地区，我们可以看到更多的阻力，其中的原因和当地土地状况的差异应不无关系。正是有着上述实例的存在，使中共中央关于查田的判断似乎有了一定的根据，但应该看到，无论是寻求宗族的保护，还是利用权力的遮蔽，地主、富农在革命大潮中寻求生存的方式都已经呈现出明显的弱者特征，保护自己的本能使他们尽一切可能苟延残存，但能够以个体的方式保存下来的仍是少数，更难以对中共革命的成效形成实质性的干扰。查田运动从蛛丝马迹中发现敌人活动的佐证，再把这些细枝末节加以放大，以此为革命的困难寻找到地主、富农破坏的因由，这是符合阶级分析的省事办法，也是多次政治运动的共同逻辑。


其实，查田运动前某些农村土地分配确实存在的问题，和社会生活中无所不在的权力因素也密切相关。江西省委发现，由于在分田中“惯用自上而下的经过红军，经过临时苏维埃政权，限期将田分好的派田方式，因此，不但政府工作人员可以特别分好田多分田，包庇分田，操纵成份的分析等等脱离群众压迫群众的现象在新发展区域发生”\footnote{《中共江西苏区省委四个月（一月至四月）工作总报告》，《江西革命历史文件汇集（1932年）》（一），第149页。}。查田中发现的所谓地主、富农多占现象更多的是和其作为干部及干部家属的身份联系在一起的。


观察查田运动，其阶级革命的性质及背景当然是首要因素，同时，支持苏维埃的财政需要也是运动开展的一个重要诱因。苏区发展早期，中共财政来源多依赖打土豪的收入，随着苏区内部打土豪的结束和苏区外围的相对固定，加上赤白对立影响和国民党方面的封锁，苏区财政供给面临严重困难。而农业收入减产，“农业税短收很大”，\footnote{《江西省各县及中心区财长联席会议（1933年10月）》，陈诚档案缩微胶卷008·661/6421/0240，中国社会科学院近代史研究所藏。}更加重了财政负担。查田运动虽然主要是一场政治运动，但其表现形式却是经济的，其中确也不难看到经济方面的考虑。运动中特别重视对被定为地主、富农者的财物没收，在6月召开的八县区以上苏维埃负责人查田运动大会上，明确要求7～9月应没收地主现款、富农捐款80万元，并出售300万元的经济建设公债。运动开展期间，临时中央政府财政和土地部门鉴于“现在红军需款很多，而各地筹款又不甚得力”，\footnote{《中央财政、土地部为筹款问题给乡主席、贫农团的一封信（1933年9月19日）》，陈诚档案缩微胶卷008·663/8847/023，中国社会科学院近代史研究所藏。}又要求各地应在运动中加紧对财物的没收、统一管理和上交。更值得注意的是，几乎与查田运动开展同时，苏区中央决定设立没收征发委员会，“专门负责管理地主罚款，富农捐款及归公没收物品等。红军初到城市向商人筹款时，也由没委会负责进行”。\footnote{《中央财政人民委员部没收征发委员会组织与工作纲要（1933年7月10日）》，《革命根据地经济史料选编》（上），江西人民出版社，1986，第451页。}没收征发委员会分三个系统组织，分别为地方财政机关系统、红军政治部系统、军区及地方武装系统。委员会在地方省、县、区均设常驻工作人员，乡委员不脱离生产，但至少须有三个人负专责，不能任其他职务。没收征发委员会“到了查田、查阶级运动得到了彻底胜利，该地方已无地主罚款、富农捐款收入”\footnote{《中央财政人民委员部没收征发委员会组织与工作纲要（1933年7月10日）》，《革命根据地经济史料选编》（上），第451页。}时，可以呈报撤销，可见该委员会实际是配合查田运动设置的专门筹款机关，查田的经济目标于此可见一斑。


查田运动开始后的7、8、9三个月，仅博生、乐安、石城、胜利四县，就利用没收财物和强制捐款“筹到了十八万元”，\footnote{《江西省各县及中心区财长联席会议（1933年10月）》，陈诚档案缩微胶卷008·661/6421/0240，中国社会科学院近代史研究所藏。}整个中央苏区共完成606916元。\footnote{《中华苏维埃共和国中央执行委员会与人民委员会对第二次全国苏维埃代表大会的报告》，《红色中华》第二次全苏大会特刊第3期，1934年1月26日。}而1932年苏维埃辖下的江西全省农业税收入只有55万元，“其中人口最多土地最多的博生县只收到八万三千元，石城全县只收二万七千元，广昌全县只收到一万五千元”，\footnote{《江西省各县及中心区财长联席会议（1933年10月）》，陈诚档案缩微胶卷008·661/6421/0240，中国社会科学院近代史研究所藏。}几个月的罚款收入几乎等于一年的农业税收入，其在支持苏维埃财政上的短期功效不可小视。福建方面收入虽比不上江西，但成绩也相当不错，福建省委报告：“省苏在纠正了非阶级路线的财政政策以后，在短时期内，在查田运动的初步开展中，罚款与筹款已达到十二万余元。”\footnote{《中共福建省委工作报告大纲（1933年10月26日）》，福建省档案馆、广东省档案馆编《闽粤赣边区革命历史档案汇编》第1辑，第385页。}而1932年福建的农业税收入是15万元。新成立的闽赣省因为处于新区、边区，可打击对象较多，被寄予的筹款希望更大，该省计划在1933年7月至次年2月通过查田向土豪筹款662000元，富农捐款208000元，\footnote{《闽赣省财政部七、八、九三个月工作和八、九、十、十一、十二、一、二七个月筹款计划》，《闽赣苏区文件资料选编》，第60～61页。}虽然由于数额过大，这一计划难以完成，但查田运动现实的筹款需求在此仍体现得非常明显。在战争环境下，这样非常规的筹款方式较之正规的财政制度便捷得多。考虑到当时巨大战争迫近，支撑战争经费难以为继，中共中央寄希望于此有其不得已而为之的可理解之处，但在稳固的苏区范围内，采用这种方式仍不免有饮鸩止渴之嫌。


对于查田运动广泛展开后暴露的问题，中共中央和毛泽东都很快有所意识。1933年底，毛泽东主持开展纠偏，一定程度上抑制了运动过火的局面。胜利县纠正了1512家错划地主、富农，而该县地主、富农总计家数为2124家，\footnote{王观澜：《胜利县继续开展查田运动经验》，《斗争》第64期，1934年5月26日。}可见错划面之广。但纠偏和当时中共中央总体思路相背离，很快被作为右倾受到批判。毛泽东以生产资料占有状况、方式和剥削量区别成分的实事求是主张被讥为“算成分”，而错划成分的不公平现象也被认为“到底是不多的，值不得我们多大的注意”。\footnote{张闻天：《关于开展查田运动中一个问题的答复》，《红色中华》第168期，1934年3月29日。}1934年春，随着反“围剿”军事的日渐紧张，查田运动再次以激烈的形式展开，人民委员会明确指示：“在暴动后查田运动前已经决定的地主与富农，不论有任何证据不得翻案。已翻案者作为无效。”“必须坚决打击以纠正过去‘左’的错误为借口，而停止查田运动的右倾机会主义。”\footnote{《关于继续开展查田运动的问题，人民委员会训令中字第一号》，《红色中华》第164期，1934年3月20日。}


显然，在中共中央看来，面对第五次反“围剿”军事不利的局面，农村的阶级关系应该更加紧张，对敌对势力的打击应更加严厉。问题是，经过数年的革命荡涤，苏区内所谓有组织的阶级敌对势力其实已更多存在于他们的想象之中，硬要追索阶级敌人的结果，只能是盲目扩大打击面，使侵害中农和乱划成分变得难以避免，正如张闻天后来所认识的：

\begin{quote}

在查田运动的名义之下，任意没收了地主富农兼商人的店铺与商品，或是把工人当了地主打，但没一个人敢起来纠正，甚至负责的机关也听其自然不去干涉，因为大家怕这种干涉会遭到右倾机会主义头衔与同地主资本家妥协的罪名。\footnote{张闻天：《反对小资产阶级的极左主义》，《斗争》第67期，1934年7月10日。}
\end{quote}


作为第五次反“围剿”前夕开展的一次社会政治运动，查田运动不失为观察当时中央苏区政治的一个良好切入点，运动的发起方式、组织推进、阶级观点、群众意志、现实功用，都具有明显的苏区政治的烙印，年轻的中共领导群体的思维路向、政治领导方式及其面临的困境，在查田运动中，似乎均可见微知著。


中央苏区是依靠军事力量建立、发展、存在的控制区域，战争的胜败始终是中央苏区兴衰的决定性因素，民众的趋向很难简单和其成败画等号，这中间的影响也无法量化，但一些不利状况的出现，终究有蚁穴溃堤的风险，应该不为任何政治力量所乐见。而就在第五次反“围剿”即将开展的关键时期，中共中央推出查田运动，本意是为反“围剿”扫清障碍、准备资源，结果却自乱阵脚，某种程度上成为寻找和制造敌人的过程，造成苏区群众的普遍恐慌，损害中共与群众间关系。红军离开后到苏区调查的学者们观察到：“查田运动开始以后，在国军第五次围剿中，匪区农民逃投国军者日多，匪兵投诚者日众。”\footnote{汪浩：《收复匪区之土地问题》，第45页。}事实上，中央苏区成规模的群众逃跑现象就是从这一时期开始的。如果说当年的中共中央对运动的恶果毫无了解，未免也太低估了他们的能力，只是有理论和现实脱节的思想基础，加之吸取资源这样的需求的催迫，很多问题的出现就是大概率事件了。

